\section{Worldsheet II: circle picture and covering}
\label{ch:worldsheet-circle}

This chapter contains the circle picture of the quartic form: the four-punctured sphere of the vertex region is cut along a circle that encloses the two punctures of one $\Sigma_2$ slot, which writes $\sprod{2}{2}$ as a gluing of two three-punctured spheres over a region of the plumbing parameter.
It expands the worldsheet subsections of the paper on the circle picture, the atlas picture, the six-point exchange and the covering question, and the single-strand route, reproduced in full as the first block, as they stood before the paper was shortened.
The sections written for this document give the derivations behind that block which are only stated there, or not at all: the quartic form on its vertex region as a dressed gluing off-shell, the measure identity that makes the rigid gluing exact on-shell, the Sonoda--Zwiebach sewing that follows from it, the coordinate data of the explicit quartic vertex, the reduction of its circle states to the rigid family, the direct tests of the gluing recipe, and the patch with three punctures.
The last two blocks come from an earlier note on the inverse of the quartic form as a circle gluing, with the block structure of the inverse in the fixed-halves model, its tower of necks and the six-point composite in the circle toy.
The dictionary between fields and punctures, the no-kernel theorem for the inverse of the quartic form and the counting of moduli on which this chapter builds are in \cref{ch:worldsheet-dictionary}.

\subsection{The circle picture in the paper}
\label{sec:wscircle:paper}

The block is the part of the paper's worldsheet section that follows the counting of moduli, in the paper's notation: $\varphi_{n,1}$ for the circle map, $\sprod{2}{2}^{\mathrm{circ}}$ for the fixed-halves model, $Q_V$ for the pre-image of the vertex region in the plumbing parameter $q = \e^{- s + \I \theta}$, $\tau_V(\theta)$ for its outer radius, $\lambda$ and $\gamma'$ for the parameters of the halves, with the cubic vertex of the theory at $\lambda = 1 / 2$, $\gamma' = \beta$.
It states the following results, in the order in which they appear.

\begin{itemize}
  \setlength{\itemsep}{0.5ex}

  \item The circle map: on an $n$-string field it produces one string state of ghost number $2 n$, with the relative momenta as level content; it is not level matched, at $n = 2$ for every momentum of a tachyon pair; the exchange of the two punctures acts as $(-1)^{L_0^-}$ up to the Grassmann sign; it is injective on inputs of finite level at $n = 2$.
    \status{mixed}\record{071-circle-map-symn-to-h (R79)}
    The circle state is exact for inputs of finite level; its extension to the infinite level sums of a string field is not established, as the block says.

  \item The fixed-halves model: the gluing over $Q_V$, the region with its outer radius, the model's inverse with its three factors, the tower of necks, the telescoping and the real zeros of the transform.
    \status{mixed}\record{074-p22-inverse-circle-gluing (R80)}
    The map from $q$ to the cross-ratio, the inverse in the model and the telescoping are proven; that $Q_V$ is star-shaped with a single-valued cutoff is numerical at $\beta = 3 / 2$, with $\tau_V(\pi)$ exact; the real zeros of $\hat f_4$, $\hat f_6$ are numerical; the dropping of the projector, and with it the reading of the zeros as poles, holds under the hypothesis on the circle states stated in the block.
    The block says that the exactness on-shell holds modulo a measure identity argued by hand; that identity is now proven (\cref{sec:wscircle:measure}), and the note inserted before the gluing paragraph says so.

  \item The convergence of the sector expansion on the fold arcs: for four tachyons, in every order of summation if and only if $\alpha > - 1$, in level order under a weaker condition, as an Abel sum always; for inputs with matter oscillators the exponent drops by at most the total level.
    \status{mixed}\record{079-fold-arc-convergence (R83); 084-fold-arc-excited}
    The every-order threshold is proven modulo the numerical shape of the region; the level-order thresholds and the Abel sum are argued from an endpoint analysis and checked numerically at six momenta.

  \item The seam mismatch: with rigid halves the gluing is not consistent with the cubic exchanges at the outer seam, by $\ln ((4 \beta^2 - 1) / 4 \beta)$ at $\theta = \pi$ at all four punctures, for every $\beta > (1 + \sqrt 2) / 2$.
    \status{mixed}\record{opencode-p22-inverse-gluing (R91)}
    This is an exclusion within the class of exactly rigid fixed-halves models: the value at $\theta = \pi$ and the independence of the first-order data from the parameters of the halves are proven, the mismatch at the eleven other seam points is numerical at $\beta = 3 / 2$, and nothing is claimed about the actual $\sprod{2}{2}$, whose coordinates move.

  \item The comparison with the theorem of Sonoda and Zwiebach: on-shell the $s$-channel exchange and the quartic form together are their sewing of the symmetric Hermitian vertex with $a = 2 \beta$ over a domain that is not a disc, and the square point lies inside $\mathcal V_{0,4}$ for $\beta > (1 + \sqrt 2) / 2$.
    \status{Proven}\record{113-offshell-dressed-gluing (R113)}
    The derivation is in \cref{sec:wscircle:sz}; the reading that the theory evades the theorem through the quartic vertex is an interpretation of these facts, not a result.

  \item The selection rule~\eqref{p:eq:ws:selection-rule}: exact for rigid halves at $\lambda = 0$, absent for $\lambda \neq 0$ and for the cubic vertex of the theory; for the quartic vertex of the explicit section the second-order coefficients vanish only on the surface $q = \pm 4 \I \beta^2$.
    \status{mixed}\record{080-dressed-ckv (R86)}
    That the rule fails on an open set of the vertex region is proven; at almost every point it holds under the hypothesis that the section is real-analytic on each tile; at every point of $Q_V$ it is proven given the formulas of the explicit section, for every $\beta \geq \tfrac{1}{2} + \tfrac{1}{\sqrt 2}$; the derivations are in \cref{sec:wscircle:coordinates}.

  \item The actual circle states: at each $q$ the circle state of the quartic vertex is that of rigid halves at a parameter $\lambda(q)$, with $\overline{\lambda(q)}$ in the antiholomorphic sector, up to an invertible rescaling of the inputs; the pointwise rank through output level two.
    \status{mixed}\record{094-ph31-actual (R92)}
    The reduction is proven given the explicit section; the pointwise rank is proven near $\abs{q} = 1$, verified by randomized exact arithmetic at almost every point, and numerical at every point of the interior; derivation in \cref{sec:wscircle:rigid}.

  \item The recipe for an explicit vertex: two copies of a fixed cubic vertex, input operators at the four legs, and the dressed plumbing operator over $Q_V$; on four matter primaries the weight $\abs{\rho(q)}^{h_1 + h_2 + h_3 + h_4}$.
    \status{mixed}\record{096-gluing-recipe (R93)}
    The recipe is proven from the definitions and the transformation law, its Möbius statements given the explicit section, and the claim that the dressing cannot be moved to the neck is argued by hand, as the block says; its theorem form, with its hypotheses, is \cref{sec:wscircle:offshell}, and its direct tests are \cref{sec:wscircle:tests}.

  \item The patches of the atlas picture: the counts, the redundancy of the two-patch description, the ghost number $4$ of the dressed patch state for every $n \geq 2$, and the generalized trees.
    \status{mixed}\record{077-npatch-atlas (R82)}
    The counts are proven; for $n \geq 3$ the statements are necessary conditions on a patch that is not constructed, and the construction at $n = 3$ is in \cref{sec:wscircle:npatch}.

  \item The six-point exchange: the exclusions of a compact region with momentum-independent boundary and of a region of sector type when the composite vertex has a pole off the plumbing lattice, the exclusion of a union of Feynman regions, and the cusp-shaped region.
    \status{mixed}\record{035-negative-moduli-composite (R39); 036-covering-m06-composite (R40)}
    The exclusion of a gluing of the inverse within the class of gluings with a Möbius cut coordinate is the result of \cref{ch:worldsheet-dictionary} on the asymmetric vertex and its gluings.
    The tests at real momenta and their calibration on the slab model:
    \status{mixed}\record{105-exchange-real-momenta (R99, R100)}
    The positivity of the measure and the bounds implied by a sub-region or by a constant weight are proven; the definition of the composite vertex as a limit of level truncations rests on the assumptions the block lists, which are open hypotheses: Gaussian damping of the quintic vertex in the Euclidean relative momentum of the pair, invertibility of the truncated pair blocks at real momentum off a discrete set, convergence of the truncations, meromorphy of the exchange on a connected set of complex momenta, and invariance of $\mathcal V_{0,6}$ under complex conjugation for the constant weight $\tfrac{5}{12}$.
    The conjecture that the composite vertex covers a sub-region of $\mathcal V_{0,6}$, the existence of a pole of the truncated composite off the plumbing lattice, and whether a cusp region reproduces it, are open.
    The level-zero kernel paragraph quotes results of \cref{ch:momentum-space}, with the status given there; the zero of the discretized determinant at complex $P^2$ is numerical.

  \item The single-strand route: the plumbing family of two copies of $\vertex{3}'$ covers the interior of $\mathcal V_{0,4}$ once for a stub length varying along a seam symmetric in the three channels, and never for a constant one.
    \status{mixed}\record{032-v3prime-route-feasibility (R37)}
    The covering is numerical and the seam under-determined; the on-shell operator identity holds under the condition on the discs and the hypothesis that the sewing converges over the reversed plumbing range, which is open; the measure identity it uses, ``argued by hand'' in the block, is the one proven in \cref{sec:wscircle:measure}.
\end{itemize}

\subsubsection{The circle picture}
\label{p:sec:ws:circle}

We look for a mechanism behind this counting.
The inverse carries two $c$ insertions and no modulus, and the identity $\vertex{4} \, \sprod{2}{2}^{-1} \, \vertex{4} = \vertex{4}$ cancels two antighosts and two moduli of its neighbors.
We now describe a model in which this is realized by cutting the four-punctured sphere of the vertex region along a circle that encloses the two punctures of one $\Sigma_2$ slot.
The model reproduces the quartic form exactly on-shell at four points, modulo a measure identity argued by hand, and everything we say about the inverse is a statement about the model.






\paragraph{Halves and circle states}
\label{p:sec:ws:halves}

Take two copies of a rigid three-punctured sphere with ordinary punctures at $z = 0, \infty$ and a special puncture at $z = 1$, whose local coordinate $w$ is
\begin{equation}
  \label{p:eq:ws:half}
  z
    = \frac{2 \beta - w}{2 \beta + w},
\end{equation}
with $\beta > 0$ the stub parameter~\cite{Erbin:2026:SL2CVerticesClosed}, and the numbers are computed at $\beta = 3 / 2$ and $2$.
In the $w$-plane the ordinary punctures sit at $w = \pm 2 \beta$.
We write $\vertex{3}'$ for this cubic vertex, asymmetric for generic coordinates at its ordinary punctures, and return to it in the outlook.
At one value of the parameters introduced below it is the cubic vertex of the theory itself.
Glue the two copies at their special punctures by $w_1 w_2 = q$, with $q = \e^{- s + \I \theta}$ the plumbing parameter, $s$ the neck length and $\theta$ the twist.
Any circle $\abs{w_1} = r$ with $\abs{q} / (2 \beta) < r < 2 \beta$ separates the two punctures of the first copy from those of the second, and a symmetric cut $r = \sqrt{\abs{q}}$ exists if and only if $\abs{q} < 4 \beta^2$.

The state that one copy puts on the circle is created by its two transported vertex operators on the vacuum,
\begin{equation}
  \label{p:eq:ws:circle-map}
  \varphi_{2,1}(A_1, A_2)
    := f_1 \circ A_1(0) \; f_2 \circ A_2(0) \, \ket{0},
\end{equation}
with $f_i$ the local coordinate maps and the product radially ordered.
Its level expansion is the repeated operator product expansion of the two operators.
The same construction with $n$ punctures inside one circle defines a map $\varphi_{n,1} \colon \Hn{n} \to \Hs$, which we call the circle map.
On an $n$-string field it produces one string state of ghost number $2 n$, whose zero mode carries the total momentum $P$ and whose level content encodes the $n - 1$ relative momenta of~\eqref{p:eq:ws:jacobi}.
The relative momenta are not lost in the merging, being recoverable from the level content, and at $n = 2$ the map is injective on inputs of finite level.
The circle state is exact for inputs of finite level, when the coordinate discs lie inside the disc of the circle.
Its extension to the infinite level sums of a string field is not established.

It is not level matched, because the configuration of punctures inside the circle is not rotation invariant.
At $n = 2$ the failure is present for every momentum of a tachyon pair, the ghost contribution to the $L_0^- = \pm 2$ components never vanishing.
With $f_2 = - f_1$, which holds for the halves above, the exchange of the two punctures acts as $(-1)^{L_0^-}$ up to the Grassmann sign, so that on a graded-symmetric $\Sigma_2$ only the even-$L_0^-$ components of the circle state survive.
The merged state enters a sewing with a closed string only through the level-matching bivector $b_0^- \delta_{L_0^-, 0}$ of the strands, which discards its components with $L_0^- \neq 0$; the part it loses is structural, not an artifact of a choice.


In their no-go theorem, Sonoda and Zwiebach allow for propagating states that are not level matched~\cite[pp.~188--189]{Sonoda:1990:CovariantClosedString}.
Their external states are level matched, but the off-shell states running in the propagator may form a larger set when the propagator is not the standard one.
Since their proof examines only amplitudes whose external states are level matched, its conclusion holds as long as the set of all off-shell states contains the level-matched ones.\footnotemark{}
\footnotetext{
  Considering different sets of the on-shell and off-shell states can be seen in different instances~\cite{Pius:2016:CutkoskyRulesSuperstring,Bautista:2019:QuantumGravityTimelike,Bautista:2020:BRSTCohomologyTimelike}.
}
The circle state belongs to such a larger set.
The gluing defined below propagates it over a region of the $q$-plane that is not a disc, so that on-shell the circle picture of the quartic form realizes the enlargement they contemplate.
We compare the circle picture with their sewing once the gluing has been defined, and we find that this enlargement is not what lets the theory evade the theorem.
Constructions of closed string field theory without level matching have been considered in~\cite{Okawa:2022:ClosedStringField,Erbin:2023:ClosedStringTheory}.

\begin{quote}
  \small
  \emph{Note added in this write-up.}
  The paragraph below, and the opening of this subsection, state the exactness of the model on-shell ``modulo a measure identity argued by hand''.
  The measure identity is now proven, by two independent routes, and the off-shell statement is stronger than the paragraph says: on its vertex region and on finite-level inputs the quartic form itself, for any vertex given as a section of local coordinates continuously differentiable on each tile (for the explicit section, under the hypothesis that it is real-analytic on each tile), is a gluing of two fixed halves over $Q_V$, with input operators and dressed antighosts, and the rigid gluing is its on-shell value.
  The derivations, and the Sonoda--Zwiebach sewing that follows from the measure identity, are in \cref{sec:wscircle:offshell}.
\end{quote}

\paragraph{Gluing over the vertex region}
\label{p:sec:ws:gluing}

With the halves kept fixed, the only deformation of the glued surface is that of $q$, a dilation or a rotation of the circle, and the antighost insertions that go with the two moduli are the zero modes in the frame of the circle,
\begin{equation}
  \label{p:eq:ws:schiffer}
  b(\pd_s)
    = - b_0^+,
  \qquad
  b(\pd_\theta)
    = \I \, b_0^-.
\end{equation}
Gluing over a region $Q$ of the $q$-plane defines the bilinear form
\begin{equation}
  \label{p:eq:ws:p-circ}
  \sprod{2}{2}^{\mathrm{circ}}(\Sigma_2, \Sigma_2')
    := - \I \kappa \int_Q \dd s \, \dd\theta \;
      \big\langle \varphi_{2,1}(\Sigma_2') \big| \, b_0^+ b_0^- \, \e^{- s L_0^+} \e^{\I \theta L_0^-} \, \big| \varphi_{2,1}(\Sigma_2) \big\rangle
\end{equation}
in the BPZ inner product, with $\kappa$ the normalization of the measure, fixed by the propagator and cancelling in every statement about the inverse.
The map from $q$ to the cross-ratio of the glued sphere is injective on $\abs{q} < 4 \beta^2$, the disc $\abs{q} \le 1$ being the $s$-channel Feynman region.
The pre-image of the vertex region is an annulus with an angle-dependent outer radius,
\begin{equation}
  \label{p:eq:ws:region}
  Q_V
    = \{\, 1 < \abs{q} < \e^{\tau_V(\theta)} \,\},
\end{equation}
mapped bijectively up to measure zero onto $\mathcal V_{0,4}$, the outer annulus $\e^{\tau_V} \le \abs{q} < 4 \beta^2$ covering the $t$- and $u$-channel Feynman regions.
That the region has this star-shaped form with a single-valued cutoff is a numerical finding at $\beta = 3 / 2$, for the vertex family of~\cite{Erbin:2026:SL2CVerticesClosed}.
The cutoff $\tau_V$ is $\pi$-periodic, with $\tau_V(\pi) = \ln \tfrac94$ in closed form, and reaches $\ln 4 \beta^2$, the largest admissible value, on two arcs where the circle passes through the ordinary punctures.
The cutoff is determined by the vertex: the one circle that $\sprod{2}{2}$ prescribes reaches all of $\mathcal V_{0,4}$, and there is no freedom of splitting the region among three channels.


Nothing collides on the two arcs.
The singularities of the integrand are the cross-channel collisions, at the two points $\theta = 0, \pi$ of the fold circle $\abs{q} = 4 \beta^2$, which lie outside the closure of $Q_V$.
They set the decay of the coefficients of the expansion of the integrand in eigenstates of $L_0^\pm$.
For four tachyons the integrand is $\abs{y}^{h_0} \abs{1 - y}^{2 a_-} \abs{1 + y}^{2 a_+}$ with $y = q / 4 \beta^2$, the exponents $a_\mp$ being set by the $t$- and $u$-channel momenta.
  The coefficient at level $h$ in a sector of fixed $L_0^-$ scales as $h^{- 2 \alpha - 2}$, with $\alpha = \min (a_-, a_+)$.
The statements below about the transform of the region, its closed forms, its real zeros and the identity it satisfies sector by sector involve no sum over sectors and do not depend on the arcs.
The arcs, where $\abs{y} = 1$, decide instead whether a matrix element of the form, an integral over $Q_V$, may be computed as a sum of sector contributions.
For four tachyons the sector decomposition converges in every order of summation if and only if $\alpha > - 1$, which on-shell is the condition that both cross channels sit below their tachyon pole, $\alpha' t, \alpha' u < - 4$.
In level order it converges under the weaker condition that both sit below their massless pole, with $\alpha' s > - 10$, by an endpoint analysis that we have checked numerically.
At every momentum the region integral is the Abel sum of the level series, by the same analysis.

For inputs with matter oscillators the exponent $\alpha$ drops by at most the total level, and by exactly that for generic polarizations.
On-shell the drop cancels against the masses, so that the every-order condition $\alpha' t, \alpha' u < - 4$ holds for all matter-oscillator inputs and is sharp for generic polarizations.
The reading of the real zeros of the transform as poles, below, is accordingly a statement about the remainder of the level-ordered sum, not about sums over sectors.

For four on-shell primaries the integrand of $\vertex{4}$ does not depend on the local coordinates.
  We have argued by hand that the measure of the gluing family is the pull-back of the canonical one.
Injectivity then gives $\sprod{2}{2}^{\mathrm{circ}} = \sprod{2}{2}$ with $Q = Q_V$: the model is exact on-shell at four points, for inputs of finite level.
On-shell the quartic form is thus the matrix element, between the two circle states, of the operator $b_0^+ b_0^- \, \e^{- s L_0^+ + \I \theta L_0^-}$ that gives the string propagator~\eqref{p:eq:ws:propagator-psi} with its twist, integrated over the region~\eqref{p:eq:ws:region} instead of the unit disc.
  The inverse of the model is explicit sector by sector, as~\eqref{p:eq:ws:inverse-circle} below: the pair $c_0^- c_0^+$ dual to the two antighosts, times the reciprocal of the transform of the region, up to a projector discussed there.

\paragraph{Seams with the Feynman regions}
\label{p:sec:ws:seams}

Off-shell this simple form cannot hold, and the reason is more than a difference of coordinates: with rigid halves the gluing is not consistent with the cubic exchanges at the boundary of its region.
The region $Q_V$ has two seams with the Feynman regions.
On the inner one, $\abs{q} = 1$, the glued surface is the $s$-channel exchange built from the same halves, and the local coordinates at the four external punctures match identically.
On the outer one, $\abs{q} = \e^{\tau_V(\theta)}$, the same surface is the $t$- or $u$-channel exchange at unit plumbing modulus, and the coordinates that the two descriptions put at the external punctures must agree up to a phase for the vertex and the exchanges to fit together off-shell, with the standard antighost insertions and level-matched external states.
They disagree already at first order.
At $\theta = \pi$ the logarithms of the mapping radii differ, at all four punctures, by
\begin{equation}
  \label{p:eq:ws:seam-mismatch}
  \ln \frac{4 \beta^2 - 1}{4 \beta},
\end{equation}
exactly, for every $\beta$ above $(1 + \sqrt 2) / 2$, the value at which this vanishes and the region pinches at $\theta = \pi$, below which the $s$- and $u$-channel regions overlap there; at the eleven other seam points we sampled, at $\beta = 3 / 2$, the mismatch is of the same size.
On the two fold arcs, an internal seam where the two sheets of the plumbing covering meet, the coordinates agree to first order and, for the halves of the cubic vertex of the theory, differ at second order by $1 / (4 \beta)$ in the second-order coefficient of the coordinate map normalized by the square of the first; they agree there only for halves whose coordinates at the ordinary punctures are pure rescalings, which those of the cubic vertex are not.
The first-order mismatch does not depend on the coordinates of the halves at their ordinary punctures, and the range is fixed by the vertex, so no gluing of rigid halves over any region reaching the outer seam, whichever member of the family of halves is used, is the quartic form off-shell, nor any quartic vertex consistent with the cubic exchanges.
About the actual $\sprod{2}{2}$, whose coordinates move, this says nothing.

This failure has the form of the obstruction in the theorem of Sonoda and Zwiebach~\cite{Sonoda:1990:CovariantClosedString}, which we meet here in a configuration next to the class their proof covers.
At the parameters of the theory, $\gamma' = \beta$ and $\lambda = 1 / 2$ in~\eqref{p:eq:ws:half-ordinary}, where the halves are two copies of its cubic vertex, the vertex at each end of the gluing belongs to the class their theorem assumes: it is their symmetric Hermitian $\group{SL}(2, \C)$ vertex, with Möbius local coordinates and the parameter $a = 2 \beta$~\cite[eq.~(2.18)]{Sonoda:1990:CovariantClosedString}.
On-shell and for inputs of finite level, the $s$-channel exchange and the quartic form together are then their sewing of that vertex over a domain that is not a disc, the union of the unit disc and $Q_V$, with propagating states that are not level matched.
One clause of their sewing hypothesis fails here, however: their sewing domain is common to the three channels, whereas here the extended domain is attached to the $s$-channel plumbing alone, the $t$- and $u$-channel exchanges keeping their unit discs.
The asymmetry has a reason that is itself outside their framework: the theory does not satisfy their assumption that the three channels cover $\mathcal M_{0,4}$ once, since the three Feynman regions of the cubic vertex leave $\mathcal V_{0,4}$ uncovered.
The extension over $Q_V$ describes that region: it is not a fourth channel of the cubic theory but the quartic vertex of the original theory.
The requirement that the coordinates agree up to a phase at the seams is their condition for the BRST invariance of the four-point amplitude~\cite[eq.~(3.49)]{Sonoda:1990:CovariantClosedString}.
Its failure here is consistent with their conclusion that no symmetric Hermitian vertex, sewn by the plumbing fixture over a sewing domain common to the three channels that makes them cover the moduli space once, has coordinates that agree up to a phase at the seams, but their proof does not imply it.
Their proof runs through the point of the moduli space where the four punctures form a square, the cross-ratio $-1$, which lies on a seam between two Feynman regions.
For the gluing of the halves above and for every $\beta$ above $(1 + \sqrt{2}) / 2$, this point and its two images $\tfrac{1}{2}$ and $2$ under the permutations of the punctures are not the image of any $q$ in the three Feynman discs, real or not: they lie inside $\mathcal V_{0,4}$, and they reach its seams only at the pinch value.
The seam between the vertex region and a Feynman region, which their construction does not have, is therefore out of reach of their argument.
We conclude that the theory evades the theorem neither through the enlargement of the space of propagating states nor through the vertex at the ends of the gluing.
The off-shell quartic form is not that gluing: its coordinate maps move with $q$ inside $Q_V$, as~\eqref{p:eq:ws:dressing} below records, and the line of $\Sigma_2$ is the inverse of the quartic form, an operator with no modulus to integrate over, which in the fixed-halves model is not a gluing, as we show below.
In their terms the four-$\Psi$ amplitude contains the elementary quartic interaction they call necessary, here written as the exchange of $\Sigma_2$, a state of two closed strings and not a closed string field, so that the coupling of $\Sigma_2$ to two string fields is not a symmetric vertex of three closed string fields.

The seam condition determines, in principle, the off-shell form.
A vertex consistent with the cubic exchanges carries, on the boundary of its region, the coordinates of the exchange diagrams.
  Over $Q_V$ the vertex prescribes its own local coordinates, which the transported ones do not reproduce: they differ at second order in the expansion around the puncture, the second-order coefficients of the actual coordinates varying over $Q_V$ while those of rigid halves are constant.
The quartic form is still a gluing over the same region, but the coordinate maps $f_i$ at the external punctures move with $q$ inside $Q_V$, and their motion dresses the two antighosts~\eqref{p:eq:ws:schiffer} by contour insertions inside the discs:
\begin{equation}
  \label{p:eq:ws:dressing}
  b(\pd_s)
    = - b_0^+ + \sum_{i} B_i[\pd_s f_i],
  \qquad
  b(\pd_\theta)
    = \I \, b_0^- + \sum_{i} B_i[\pd_\theta f_i],
\end{equation}
with $B_i[\pd f_i]$ the antighost contour inside the $i$-th disc weighted by the vector field that the variation of $f_i$ induces: the count of two antighosts is kept, their zero-mode form is not.
For the explicit section of~\cite{Erbin:2026:SL2CVerticesClosed} the motion has a simple description, which we give below once the family of halves is parametrized: at each $q$ the circle state of the quartic vertex is that of a rigid half with a parameter that varies with $q$, up to a rescaling of the inputs, so that pointwise questions about the actual form reduce to the rigid family.
After that description we give the construction behind it, a four-step recipe which applies to any explicit vertex, Möbius or not, and produces the halves, the input operators and the dressed plumbing operator.
We do not obtain the inverse of the actual form, which depends on the form on all of $\Hn{2}$, in closed form.
  We state what is known about it after~\eqref{p:eq:ws:inverse-circle}.
Everything explicit below is derived for the \emph{fixed-halves model} $\sprod{2}{2}^{\mathrm{circ}}$ with $Q = Q_V$, and the inverse and its singularities are properties of the model.

\paragraph{Inverse in the fixed-halves model}
\label{p:sec:ws:model-inverse}

In the model the inverse is not a gluing.
\begin{quote}
  \small
  \emph{Note added in this write-up.}
  ``Not a gluing'' is meant in the sense the paragraph derives: the inverse is not a gluing over any region of the $q$-plane with the standard measure and antighost insertions, and, by the ghost-number count of the no-kernel theorem, not a kernel of the standard sewing class.
  It is not the statement that the inverse cannot be built from two surfaces: the tower of necks below is a signed sum of cylinders with $\delta'$ weights.
  The model itself equals the quartic form on-shell, the measure identity being proven (\cref{sec:wscircle:measure}); every off-shell statement about its inverse is a statement about the model.
\end{quote}
In the eigenbasis $\ket{h, m}$ of $(L_0^+, L_0^-)$ the operator of~\eqref{p:eq:ws:p-circ} is $G = - \I \kappa \, b_0^+ b_0^- \hat F$ with $\hat F$ diagonal, $\hat F \ket{h, m} = \hat f_m(h) \ket{h, m}$, where
\begin{equation}
  \label{p:eq:ws:transform}
  \hat f_m(h)
    := \int_{Q_V} \dd s \, \dd\theta \; \e^{- s h + \I m \theta}
\end{equation}
is the Laplace transform in the neck length and the Fourier transform in the twist of the region.
The form lives on the block $W := c_0^+ c_0^- \Hs^\downarrow$, with $\Hs^\downarrow := \ker b_0^+ \cap \ker b_0^-$, and is non-degenerate wherever $\hat f_m(h) \neq 0$.
Three factors make up the inverse.
The first is the pair of ghost zero modes dual to the two antighosts~\eqref{p:eq:ws:schiffer}, since $b_0^+ b_0^- \, c_0^- c_0^+ = 1$ on $\Hs^\downarrow$.
They are dual to the conformal Killing vectors of the circle, a dilation and a rotation, and they are the pair of $c$'s that the counting requires.
The second is the reciprocal $1 / \hat f_m(h)$ of the region's transform.
The third is the projector onto the states the circle can produce, which enters because $\sprod{2}{2}^{-1}$ is defined on $\Hn{2}$ and pushed to $\Hs$ through the circle map, and which is invisible between two circle states.
Together,
\begin{equation}
  \label{p:eq:ws:inverse-circle}
  \varphi_{2,1} \, \big( \sprod{2}{2}^{\mathrm{circ}} \big)^{-1} \, \transp{\varphi_{2,1}}
    = P_{\mathcal I} \; \frac{\I}{\kappa} \; c_0^- c_0^+ \; \hat F^{-1},
\end{equation}
with $\transp{\varphi_{2,1}}$ the transpose of the circle map for the pairings and $P_{\mathcal I}$ the projector onto the image of the circle states.

\paragraph{Rank of the circle map}
\label{p:sec:ws:rank}


The statements about poles below rest on the question whether the projector acts as the identity on the whole even-$L_0^-$, ghost-number-four part of the block at fixed total momentum, modulo the kernel of $G$.
  The matter part of the map is triangular in level with an invertible leading term, so we expect the answer to be decided in the ghost sector, as it is, by a symmetry of the halves.
The halves~\eqref{p:eq:ws:half} carry at their ordinary punctures local coordinates of Möbius form,
\begin{equation}
  \label{p:eq:ws:half-ordinary}
  z
    = \frac{v}{\gamma' + \lambda v}
\end{equation}
at $z = 0$ in terms of a local coordinate $v$, and its mirror image at $z = \infty$, with a scale $\gamma'$ and a parameter $\lambda$.
Unless another value is stated, the numbers quoted in this section were computed at $\lambda = 0$, where the two maps are pure rescalings and the circle states obey the exact selection rule derived next.
The cubic vertex of the theory sits at $\lambda = 1 / 2$, as explained below.
The region $Q_V$ and the transform do not depend on $\lambda$, only the circle states do.
At $\lambda = 0$ the dilation of the sphere that fixes the two ordinary punctures reads exactly $+ v \pd_v$ in one ordinary frame and $- v \pd_v$ in the other.
In the coordinate $u$ of the circle state, in which the punctures sit at $u = \pm u_0$ with $u_0 = 1 / (2 \beta)$, it is the conformal Killing vector $\beta (u^2 - u_0^2) \pd_u$.
The conservation law of its antighost and stress-tensor contours, deformed from the circle onto the two coordinate discs, picks up $b_0^-$ and $L_0^-$ on the inputs, which annihilate level-matched states.
There remains a pair of identities on the output, valid for inputs of every level:
\begin{equation}
  \label{p:eq:ws:selection-rule}
  \big( b_1^- - u_0^2 \, b_{-1}^- \big) \, \varphi_{2,1}(\Sigma_2)
    = 0,
  \qquad
  \big( L_1^- - u_0^2 \, L_{-1}^- \big) \, \varphi_{2,1}(\Sigma_2)
    = 0,
\end{equation}
with $b_n^- := b_n - \bar b_n$ and $L_n^- := L_n - \bar L_n$.
The circle states therefore lie in a proper subspace of the even block, and from output level two on the level-projected circle map misses part of the block.
With the matter truncated to two flat components, it misses $5$ of the $73$ directions of the block through output level two and $60$ of $487$ through level three, and the two identities together account for the missing directions exactly.
The count through level two is the same with three components, while the count through level three depends on the matter content.
In the model at $\lambda = 0$ the projector is not the identity, however large the level of the inputs.

The rule is a property of that vertex, not of the quartic form.
For $\lambda \neq 0$ the vector field acquires a $v^2 \pd_v$ term in the local frames, the contours pick up $b_1^-$ and $L_1^-$ on the inputs, which do not annihilate them, and no selection rule follows.
The cubic vertex of the theory is a member of the family with $\lambda \neq 0$.
For $\gamma' = \beta$ and $\lambda = 1 / 2$, the coordinate~\eqref{p:eq:ws:half-ordinary} at $z = 0$, its mirror image at $z = \infty$ and the coordinate~\eqref{p:eq:ws:half} at $z = 1$ are the three local coordinates of the cubic vertex of~\cite{Erbin:2026:SL2CVerticesClosed} in the frame where its punctures sit at $0$, $1$ and $\infty$, so that at these values the halves are two copies of it.
The value $1 / 2$ is determined by the symmetry $z \mapsto z / (z - 1)$ of the vertex, which fixes the puncture at $z = 0$ and exchanges the other two.
In any frame with two of the punctures at $0$ and $\infty$, the coordinates there take the form~\eqref{p:eq:ws:half-ordinary}, in the chart $1 / z$ at $\infty$, with parameters $\lambda_0$ and $\lambda_\infty$ that obey $\lambda_0 \lambda_\infty = 1 / 4$, so that no such frame brings the cubic vertex to $\lambda = 0$.
Equivalently, the dilation $z \pd_z$ reads $\pm (v + v^2 / (2 \beta)) \pd_v$ at both punctures, never $\pm v \pd_v$, and the selection rule does not hold for the cubic vertex of the theory.
At $\lambda = 1 / 2$, with the matter truncated to two flat components and also with three, the map is onto the block through level three for inputs of level up to two, at every $\gamma'$, which only rescales $v$.
Since the entries are rational in $\lambda$, the same holds at all but finitely many $\lambda$.

For the quartic vertex itself the local coordinates move with $q$ inside $Q_V$.
Its symmetry under relabeling of the punctures is incompatible with coordinates that are dilations about the pair at every $q$.
The rule therefore fails at least on an open set of the vertex region for inputs with ghost excitations.
Analytic dependence on the modulus would carry the failure to almost every point, but the coordinates of the quartic vertex of~\cite{Erbin:2026:SL2CVerticesClosed} are at most piecewise real-analytic: the vertex region is tiled by the twelve images of a fundamental domain under relabeling of the punctures and complex conjugation, and the second-order coefficients are once but not twice continuously differentiable across six arcs where tiles meet.
Assuming, as the explicit section indicates, that the coordinates are real-analytic on each tile, the failure reaches almost every point: on the boundary with the $s$-channel region the coordinates are those of the glued cubic vertices, whose poles sit at none of the punctures, continuity carries this into the tile that touches that boundary, and the symmetry under relabeling and complex conjugation carries it to the other eleven.
For the explicit section of~\cite{Erbin:2026:SL2CVerticesClosed} the second-order coefficients of the two coordinates of a pair vanish at no point of the vertex region other than a single surface, at $q = \pm 4 \I \beta^2$ on its boundary, where the symmetry under the exchange of the pairs requires them to vanish.
For that section the rule fails everywhere else in the region, with no assumption on analyticity.
On inputs whose ghost part is $c_1 \bar c_1$ the first identity holds for every choice of coordinates and every $q$.
Over the $s$-channel Feynman region the coordinates are those of two glued copies of the cubic vertex of the theory, and on the boundary $\abs{q} = 1$ of the vertex region the coordinates of the quartic vertex agree with them.
Since the cubic vertex sits at $\lambda = 1 / 2$, the rule does not hold there either.
The cubic vertex does not have the dilation coordinates on which the rule rests, and neither does the quartic vertex, on the boundary $\abs{q} = 1$ of its region and, for the explicit section, at every point of the region but the one surface above.
The value $\lambda = 0$ is therefore not representative of the quartic form off-shell, and the selection rule is a property of rigid halves at that value.
For questions about the rank of the circle map we use the model at $\lambda = 1 / 2$, the value of the cubic vertex of the theory, whose constant second-order coefficient stands in for the varying one of the quartic vertex.
For the explicit section this is the motion of the coordinates announced above: at each $q$ in the vertex region the circle state of the quartic vertex is that of the rigid halves~\eqref{p:eq:ws:half-ordinary} at $\gamma' = \beta$ and a parameter $\lambda(q)$, with $\overline{\lambda(q)}$ in the antiholomorphic sector, composed with an invertible rescaling of the inputs, diagonal in the level, so that pointwise questions about the actual circle map reduce to the rigid family.
The parameter $\lambda(q)$ equals $1 / 2$ on the boundary $\abs{q} = 1$, vanishes on the single surface $q = \pm 4 \I \beta^2$ above and nowhere else in the region.
The rank of the pointwise map follows: through output level two, for inputs of level up to two with two flat matter components, at $\beta = 3 / 2$ and generic total momentum, the map at $q$ is onto the block near $\abs{q} = 1$ and at almost every point of the vertex region, numerically at every point of its interior, and at the surface $q = \pm 4 \I \beta^2$ it misses $5$ of the $73$ directions.
This pointwise statement is necessary for the inverse of the integrated form, and not sufficient.
The model at $\lambda = 0$ remains the case in which the selection rule holds exactly, and the numbers computed there are statements about that model.
On-shell the local coordinates drop out and the model is exact for every $\lambda$.

\paragraph{Recipe for an explicit vertex}
\label{p:sec:ws:recipe}

For a given explicit quartic vertex, Möbius or not, the asymmetric cubic vertex, the input operators and the plumbing operator off-shell are obtained as follows.
\begin{enumerate}
  \setlength{\itemsep}{0.5ex}

  \item \emph{Read the vertex in the frame of the halves.}
    The vertex is a section of local coordinates: at each point $Z$ of $\mathcal V_{0, 4}$, four maps $f_i(\cdot \,; Z)$ from the unit disc to the sphere with punctures at $0$, $1$, $\infty$ and $Z$.
    The glued sphere at $q$ carries the coordinate $u$ of the circle state, with the punctures of the first copy at $\pm u_0$ and those of the second at $\pm 2 \beta / q$, and its cross-ratio $Z(q)$ runs over $\mathcal V_{0, 4}$ as $q$ runs over $Q_V$.
    The Möbius map that sends the first three punctures to $0$, $1$, $\infty$ is unique and sends the fourth to $Z(q)$.
    Its inverse transports the four coordinates of the vertex at $Z(q)$ to germs $F_i(\cdot \,; q)$ at the four punctures of the glued sphere.
    Nothing is chosen at this step: the frame is fixed by the special coordinate $w$ of the half, the germs by the vertex.

  \item \emph{Split each germ into a fixed coordinate of the half and an input operator.}
    Choose at each ordinary puncture a coordinate $h_i$ independent of $q$ that sends $0$ to the puncture, and write $F_i = h_i \circ g_i$ with $g_i(q)$ a germ fixing $0$.
    By the transformation law of a vertex operator under a change of local coordinate, $F_i \circ A_i(0) = h_i \circ (U_{g_i} A_i)(0)$, with $U_{g_i}$ the exponential of the $L_{n \geq 0}$ and $\bar L_{n \geq 0}$ that represents $g_i$, so that the state on the special leg at $q$ is the circle state of the fixed half with coordinates $h_1, h_2$ on the inputs $U_{g_1} A_1$ and $U_{g_2} A_2$.
    The split is exact for every admissible choice of the $h_i$, and this choice, together with the placement of the inner antighost contours discussed below, is the freedom of the construction: the glued surface and its region, the measure $\dd s \, \dd\theta$ and the total antighost insertion of each tangent vector are fixed by the vertex; the coordinates of the halves, and with them the input operators, are not.
    The canonical choice is the coordinates of the cubic vertex of the theory, $\lambda = 1 / 2$ and $\gamma' = \beta$ in~\eqref{p:eq:ws:half-ordinary}.
    For a section whose coordinates on $\abs{q} = 1$ are those of the $s$-channel exchange, the input operators are trivial on level-matched inputs at that seam if and only if the $h_i$ are these, up to a phase rotation.
    The half of the $\Sigma_2$ exchange is then $\vertex{3}$ itself, and the plumbing operator reduces at the seam to the undressed $b_0^+ b_0^-$ of the model once the derivatives of the germs vanish there as well, as they do for the explicit section.
    This choice is canonical but not required: the surface states are the same for every split, and another choice of the $h_i$ moves a $q$-independent factor into the plumbing operator.

    For the section of~\cite{Erbin:2026:SL2CVerticesClosed} the germs are Möbius, and relative to the canonical half $g_i(W) = \rho W / (1 + \nu W)$ with $\nu = \rho \, (\lambda(q) - 1 / 2) / \beta$, so that
    $U_{g_i} = \rho^{L_0} \bar\rho^{\bar L_0} \, \e^{- \nu L_1 - \bar\nu \bar L_1}$,
    where $\lambda(q)$ is the parameter above, read off the position of the pole of the germ, and $\rho(q) = F_i'(0; q) / h_i'(0)$ is the mapping radius of the coordinate of the vertex relative to that of the half.
    One pair $(\lambda, \rho)(q)$ serves the four legs, by the symmetry of the section under relabeling of the punctures.
    Taking instead the half at $q$ to be the member $\lambda(q)$ of the family, with the rescaling alone as input operator, is the description given above.
    The two are the same surface states with the same total antighost insertion, and only the fixed half extends to sections that are not Möbius.
    For such a section, minimal area or known numerically, $U_{g_i}$ is $\rho^{L_0} \bar\rho^{\bar L_0}$ times an exponential of the $L_{n \geq 1}$ and $\bar L_{n \geq 1}$ whose coefficients follow order by order from the Taylor coefficients of $g_i / \rho$, the first few of which are tabulated in~\cite{LeClair:1989:StringFieldTheory-1,Rastelli:2001:TachyonPotentialsStar}.
    On inputs of level at most $L$ only the $L_n$ with $n \leq L$ act, so the $(L + 1)$-jet of the germ at $q$ determines $U_{g_i}$ there exactly.
    A numerical section has to supply this jet, with its first derivatives in $s$ and $\theta$.

  \item \emph{Assemble the plumbing operator.}
    The contribution of the vertex region, in the coordinate $q$, is the integral of the form of the section on the tangent vectors $\pd_s$ and $\pd_\theta$, with the four inputs inserted through the germs $F_i$ and the antighosts~\eqref{p:eq:ws:dressing}:
    \begin{multline}
      \label{p:eq:ws:recipe-plumbing}
      \sprod{2}{2}(A_1 \otimes A_2, A_3 \otimes A_4)
        = \kappa \int_{Q_V} \dd s \, \dd\theta \;
      \\
          \big\langle \varphi_{2,1}(U_{g_3} A_3, U_{g_4} A_4) \big| \, b(\pd_s) \, b(\pd_\theta) \, \e^{- s L_0^+} \e^{\I \theta L_0^-} \, \big| \varphi_{2,1}(U_{g_1} A_1, U_{g_2} A_2) \big\rangle,
    \end{multline}
    with $\varphi_{2,1}$ the circle map of the fixed half and the same normalization as in~\eqref{p:eq:ws:p-circ}.
    The dressing $B_i[\pd_s f_i]$ acts on its input as $A_i \mapsto U_{g_i} \, b^{(i)}[v_i] \, A_i$, where $b^{(i)}[v_i]$ is the antighost contour in the frame of the input weighted by the vector field $v_i(W) = \pd_s g_i(W) / g_i'(W)$ of the variation, and the input operator is applied after it.
    Since $v_i$ vanishes at the puncture, the contour contains the modes $b_k$ with $k \geq 0$ only, and on inputs of level at most $L$ only those with $k \leq L$.
    For the Möbius section the dressing on level-matched inputs is explicit, $B_i[\pd_s f_i] = (\pd_s \ln \abs{\rho}) \, b_0^{+ (i)} - (\rho / \beta) \, \pd_s \lambda \; b_1^{(i)} - (\bar\rho / \beta) \, \pd_s \bar\lambda \; \bar b_1^{(i)}$, and the same with $\pd_\theta$: the extra modes are $b_1$ and $\bar b_1$ on each input, with coefficients that are real-analytic on each tile of the vertex region, on the assumption stated above, and not smooth across the arcs where tiles meet.
    At the outer seam the logarithm of $\abs{\rho}$ is minus the mismatch~\eqref{p:eq:ws:seam-mismatch}, numerically at $\theta = \pi$ for $\beta = 3 / 2$ and $2$: the input operators carry the first-order mismatch that rigid halves miss.

    The dressing cannot be moved from the inputs to the neck as a change of the special coordinate: the surface states with coordinates $(h_1 \circ g_1, h_2 \circ g_2, w)$ and $(h_1, h_2, w \circ x)$ coincide only if a Möbius map of the sphere fixing the three punctures relates the two triples, and that map is the identity.
    On level-matched inputs the phase rotations of the ordinary coordinates are the only further identifications, and they act trivially.
    Level by level it can: the conservation law of the half for a polynomial vector field with a pole at the special puncture, whose jets at the two ordinary punctures match those of $v_1$ and $v_2$, trades the two inner contours, and by exponentiation the two input operators, for an insertion on the special leg, which on inputs of level at most $L$ carries creation modes $b_{-n}$ and $L_{-n}$ of the special coordinate with $n$ up to $2 L + 2$ and coefficients depending on $q$.
    We have argued this and not checked it.
    Whether a level-uniform operator on the special leg exists, other than a change of its coordinate, is open.

  \item \emph{Evaluate on primaries.}
    On inputs $A_i = c_1 \bar c_1 V_i$ with $V_i$ a matter primary of any momentum, every mode $b_{k \geq 0}$ and every $L_{n \geq 1}$ annihilates $A_i$, so the dressings drop out and the input operators reduce to the weights $\abs{\rho}^{h_i}$, with $h_i$ the eigenvalue of $L_0^+$ on $A_i$, $\alpha' k_i^2 / 2 - 2$ for a tachyon:
    \begin{multline}
      \label{p:eq:ws:recipe-primaries}
      \sprod{2}{2}(A_1 \otimes A_2, A_3 \otimes A_4)
        = - \I \kappa \int_{Q_V} \dd s \, \dd\theta \; \abs{\rho(q)}^{h_1 + h_2 + h_3 + h_4}
      \\
          \times \big\langle \varphi_{2,1}(A_3, A_4) \big| \, b_0^+ b_0^- \, \e^{- s L_0^+} \e^{\I \theta L_0^-} \, \big| \varphi_{2,1}(A_1, A_2) \big\rangle,
    \end{multline}
    the integrand of the model, which on such inputs does not depend on $\lambda$, times one function of $q$.
    For a section that is not Möbius the weight is $\prod_i \abs{\rho_i(q)}^{h_i}$ with the four mapping radii $\rho_i = F_i'(0; q) / h_i'(0)$, which need not be equal.
    The whole off-shell content of the vertex on such inputs is thus the mapping radii of its coordinates relative to those of the half, as a function of $q$.
    On-shell the weights are one, and the formula is the exactness stated above.
\end{enumerate}
We have checked the recipe numerically on the section of~\cite{Erbin:2026:SL2CVerticesClosed}.
At $\beta = 3 / 2$, on inputs that are not primary, with matter descendants at complex momentum and ghost excitations, the input operators reproduce the four-point surface states of the vertex to machine precision at three points of the vertex region, and the sum over intermediate states converges to them at the two points with $\abs{y} = 0.12$ and $0.22$, $y = q / 4 \beta^2$.
At $\abs{y} = 0.36$ the truncated sum is consistent with them but not conclusive.
On the zero-momentum tachyon $c_1 \bar c_1 \ket{0}$ and ghost dilaton, the dressed integrand of~\eqref{p:eq:ws:recipe-plumbing} agrees with the integrand of the vertex on these states in~\cite{Erbin:2026:SL2CVerticesClosed}, to $10^{-9}$, at three points of the region for $\beta = 3 / 2$ and one for $\beta = 2$.
These states are annihilated by the $b_0^{(i)}$, whose coefficients in the dressing are therefore not tested.

Three statements of this section rest on the projector acting as the identity: the inverse~\eqref{p:eq:ws:inverse-circle} with the projector dropped, the reading of the real zeros of the transform as poles, and the description of the six-point exchange as a signed density, given in the six-point paragraph below.
In the model at $\lambda = 1 / 2$ they hold under that hypothesis, under which $\sprod{2}{2}^{-1}$ acts on fields of unbounded level.
The hypothesis has finite-level support there, as at all but finitely many $\lambda$.
In the model at $\lambda = 0$ they are replaced by their projected forms, stated at each place, and the reading of the zeros as poles becomes an open criterion.
The identification of the two $c$'s with the zero modes of the circle frame is a fixed-halves statement: for $\sprod{2}{2}$ itself the dressed antighosts~\eqref{p:eq:ws:dressing} contain modes $b_n$, $n \neq 0$, and only the ghost-number count transfers.
The circle states of $\sprod{2}{2}$ depend on $q$ as well, so that the form does not factor through a single circle map and a single gluing operator as in~\eqref{p:eq:ws:p-circ}.
Once the circle map is allowed to vary with $q$, the part of this linear algebra that uses only the factorization of the form through the circle states, the inverse on them with its projector and the identity $G \Pi G = G$ there, with $\Pi := \varphi_{2,1} \, (\sprod{2}{2}^{\mathrm{circ}})^{-1} \, \transp{\varphi_{2,1}}$ the left-hand side of~\eqref{p:eq:ws:inverse-circle}, holds for $\sprod{2}{2}$ itself with $G$ the dressed gluing operator.
The block form of the inverse does not, since the dressed operator has no block structure: the inverse of $\sprod{2}{2}$ off-shell is not obtained in closed form by this route, and whether a closed form of another shape exists is open.
The pointwise rank of the actual circle map stated above is what we know, and it does not decide this.

\paragraph{Tower of necks}
\label{p:sec:ws:necks}


The reciprocal of the transform is a tower of necks.
For a rotation-independent cutoff $\tau$, which the actual region does not have, the twist integral gives $2 \pi \delta_{m, 0}$ and
\begin{equation}
  \label{p:eq:ws:tower}
  \hat f_0(h)
    = 2 \pi \, \frac{\e^{\tau h} - 1}{h},
  \qquad
  \frac{1}{\hat f_0(h)}
    = \frac{h}{2 \pi} \sum_{n \ge 1} \e^{- n \tau h}
  \quad (\Re h > 0),
\end{equation}
the Laplace transform of $\sum_{n \ge 1} \delta'(s - n \tau)$: a discrete tower of necks of lengths $n \tau$, each carrying one derivative in its length, in operator form $c_0^- c_0^+ L_0^+ \e^{- n \tau L_0^+}$.
This is the zero-mode part of the kinetic operator $c_0^- Q_B$ on a cylinder of length $n \tau$, and the derivative acts on whatever is glued to the cylinder.
It is not the integral over a reversed region: the transform of a slab and that of its mirror image multiply to $(\e^{\tau h} - 1)(1 - \e^{- \tau h}) / h^2 \neq 1$, and only a point inverts to a point.
Up to $c_0^- c_0^+$ in place of $Q_B$, the tower is the kinetic operator of an auxiliary field with one string leg obtained by gluing over the reversed plumbing range, $\abs{q} > 1$, that is, necks of negative length $s$, the range in which the region~\eqref{p:eq:ws:region} lies.
Between the $2$-string formulation and a single-string one the roles of region and tower are exchanged.
In the level-matched sector the reciprocal is regular at $h = 0$, so that the $\Sigma_2$ line has no mass-shell pole, as for the open-string auxiliary of~\cite{Erbin:2024:OpenStringStub}.
Its poles sit at complex momenta.
For the actual cutoff $\tau_V(\theta)$ the odd Fourier coefficients vanish, consistently with the circle states, and the level-matched reciprocal is again a tower, on the reversed branch $\Re h < 0$.
The transform $\hat f_0$ has no real zero, so the reciprocal has no real pole.
In the even sectors $m \neq 0$, however, $\hat f_m$ is real for real $h$, with asymptotic signs fixed by the geometry of the arcs and of the minima of $\tau_V$, and $\hat f_4$ has an odd number of real zeros.
Numerically at $\beta = 3 / 2$ it has one, at $h_* \approx - 2.36$, and $\hat f_6$ has two.
These are zeros of the block operator at real timelike momentum through the circle, $\tfrac{\alpha'}{2} (- P^2) = N + \bar N - 2 - h_*$, in sectors that contain no level-matched state and are never reached at real Euclidean momentum.
In the model at $\lambda = 1 / 2$, under the hypothesis on the circle states, they are poles of $\sprod{2}{2}^{-1}$, a statement about the remainder of the level-ordered sum, as noted above.
At $\lambda = 0$ the inverse is the compression of $\hat F$ to the circle states, and its poles are the zeros of the determinant of that compression, in a level truncation.
A real zero of $\hat f_4$ or $\hat f_6$ is one of them if the circle states contain a vector supported in that single sector, and generically only then.
Otherwise the zero moves to a momentum that depends on the compression, possibly again real and timelike, rather than disappearing.
We do not know whether the circle states contain such a vector, nor where the moved zeros sit.
The counts through level three do not reach the sectors $m = 4, 6$.
Their fate in the composite vertex at six points is discussed in \cref{p:sec:ws:six-point}.

\paragraph{Telescoping and trinion couplings}
\label{p:sec:ws:telescoping}

The identity $\vertex{4} \, \sprod{2}{2}^{-1} \, \vertex{4} = \vertex{4}$ is now a telescoping.
With $\vertex{4}(\Psi^2, \Sigma_2) = \sprod{2}{2}(\Psi \otimes \Psi, \Sigma_2)$ the chain is $\langle \varphi_{2,1}(\Psi \Psi) | G \Pi G | \varphi_{2,1}(\Psi \Psi) \rangle$, and $G \Pi G = G$ on the circle states settles it at the operator level.
Sector by sector on the block,
\begin{equation}
  \label{p:eq:ws:identity-sectors}
  b_0^+ b_0^- \, \hat f \;\cdot\; c_0^- c_0^+ \, \hat f^{-1} \;\cdot\; b_0^+ b_0^- \, \hat f
    = b_0^+ b_0^- \, \hat f,
\end{equation}
which realizes the counting: the $c_0^\pm$ cancel the $b_0^\pm$ of one neighbor's gluing, and one factor $\hat f$, one two-dimensional modulus integral as a number per sector, is cancelled by $\hat f^{-1}$.
In the neck length, for a constant cutoff, the product $\hat f \hat f^{-1} \hat f$ is the convolution of the indicator $\mathbf 1_\tau$ of the slab $[- \tau, 0]$ with the tower and with $\mathbf 1_\tau$ again,
\begin{equation}
  \label{p:eq:ws:telescoping}
  \mathbf 1_\tau * \Big( \sum_{n \ge 1} \delta'_{n \tau} \Big) * \mathbf 1_\tau
    = \sum_{n \ge 1} \pd_S (\mathbf 1_\tau * \mathbf 1_\tau)(S - n \tau)
    = \mathbf 1_\tau(S),
\end{equation}
with $S$ the total neck length: $\mathbf 1_\tau * \mathbf 1_\tau$ is a tent of base $2 \tau$, its derivative is $+ 1$ on one half and $- 1$ on the other, and the shifted copies cancel slab by slab, leaving the original slab.
The $c$'s differentiate the convolved neck modulus of the neighbors, and the composition is an alternating sum of shifted slabs rather than a product of regions.
For the actual cutoff the identity holds sector by sector as an identity of numbers, the slab picture existing only on the reversed branch.

The trinion couplings of the cubic action have the same reading.
Pointwise in the moduli of a vertex region, the integrand of a coupling of three fields $\Sigma_{k_1}, \Sigma_{k_2}, \Sigma_{k_3}$ factorizes as a three-punctured-sphere vertex composed with the circle states of the three clusters, each dressed by the antighosts of its cluster.
This holds surface by surface, with the choice of the three circles over the region as the non-canonical part, and does not give a single rigid three-vertex over the whole region.
In this sense the couplings of the cubic action are cubic vertices on circle states, at the level of the integrand.
The picture is not a map between the cubic action and a single-string theory: the circle state of $\Sigma_2$ is not the auxiliary field of the single-strand route of the outlook.
The two quadratic forms differ off-shell for a constant stub length.
For a stub length varying along a seam symmetric in the three channels, a family we describe in the outlook at the end of this section, they agree on-shell under conditions stated there.
We expect, without having shown it, that they differ off-shell.

The zero-mode identities, the closed forms and the telescoping are exact and have been checked numerically, and the real zeros of $\hat f_{4,6}$ are found numerically, the existence of one for $m = 4$ following from the asymptotics.
Under the hypothesis on the circle states the model has poles at real timelike momentum, so the hypothesis is not harmless.
None of this contradicts the exclusion~\eqref{p:eq:ws:ghost-obstruction}: the tower is infinite and discrete, it is the auxiliary field's own kinetic operator, and the $c$'s are not insertions on a surface but the duals of the neighbors' antighosts.

\subsubsection{Patches in the atlas picture}
\label{p:sec:ws:atlas}

The halves of the circle picture are the first case of a general construction, and the counting of moduli determines what the general case must contain.
In the atlas picture of closed string field theory, a sphere with $N$ punctures in a Feynman region is covered by $N$ puncture discs and $N - 2$ three-punctured spheres as charts.
The charts are joined by $N - 3$ transition functions $w_a w_b = q_{ab}$ whose parameters are the moduli, each carrying the two antighosts of its dilation and rotation.
At genus zero a global coordinate exists, and the three-punctured spheres are dispensable as charts but not as vertices: their local coordinates are what fixes the coordinates at the punctures.
A patch containing $n$ punctures is, as a surface, an $(n + 1)$-punctured sphere with a distinguished leg, the boundary of the chart, on which the state read is the circle state $\varphi_{n,1}$ of the $n$ inserted states.
Two such patches joined by one transition function give the $2 n$-punctured sphere of the kinetic term of $\Sigma_n$.
  This description is redundant by four complex parameters: the Möbius freedom of the global coordinate and the rescaling of one chart coordinate against $q$.
The dilation that fixes the two chart centers is the seam: it moves the collar circle and re-splits the surface between the two patches without changing it.
The two circle states change by opposite dilations that cancel in the pairing.

The counts decide what a patch must contain.
Two rigid patches and one transition function give one complex modulus, against the $2 n - 3$ of the $2 n$-punctured sphere, and the two agree only at $n = 2$.
There the patches are the halves above, and the region of the transition function is the region~\eqref{p:eq:ws:region}, which covers $\mathcal V_{0,4}$ once with the qualifications stated there.
For $n \geq 3$ each patch must carry $n - 2$ internal moduli, the full moduli of an $(n + 1)$-punctured sphere.
The patch then has the moduli and antighosts of the chart of a higher vertex with one distinguished leg: an asymmetric vertex region with its own local coordinates.
We do not know whether the symmetric vertex region, with the coordinate of one leg made special, is suitable.
The region of the transition function must be fibered over the internal moduli so that the three together cover the vertex region of the kinetic sphere once.
This is a necessary condition on an object we have not constructed.

The dressed state on the boundary of a patch, the circle state with the antighosts of the internal moduli, has ghost number $2 n - 2 (n - 2) = 4$ for every $n \geq 2$, and two of them joined by a transition function give the six of the sphere.
The trinion couplings of the auxiliary fields are a rigid three-punctured sphere with three patches attached, each patch leg contributing $n - 1$ complex moduli, and the count closes for every triple of cluster sizes.
The resulting tree has two kinds of edges.
The gluing edges of the transition functions each add one complex modulus and two antighosts.
The contraction edges of the $\Sigma_k$ lines each remove the $2 k - 3$ complex moduli of the kinetic sphere they sit on, and trade its $4 k - 6$ antighosts for as many ghosts.
For any such tree the number of moduli is the number of punctures minus three, as for the ordinary trivalent trees, so counting cannot distinguish the two expansions.
The difference lies in the regions and in the local coordinates.

\subsubsection{The six-point exchange and the covering question}
\label{p:sec:ws:six-point}


The composite vertex $\vertex{5 \cdot 5}$ enters the reduced sextic vertex~\eqref{p:eq:hs6:dressed-vertex}.
  The six-point amplitude is the same integral over all of $\mathcal M_{0,6}$ in both theories.
On one side it is the sum of the Feynman regions of the string field and of the vertex region $\mathcal V_{0,6}$~\cite{Zwiebach:1993:ClosedStringField}; on the other, the same Feynman regions, the reduced vertex $\dvertex{6}$ and $\tfrac{6}{5} \vertex{5 \cdot 5}$.
We conjecture that the two sides agree region by region on $\mathcal V_{0,6}$: $\tfrac{6}{5} \vertex{5 \cdot 5}$ is the integral of the standard measure over a sub-region of $\mathcal V_{0,6}$ and $\dvertex{6}$ the integral over its complement, so that together they cover $\mathcal V_{0,6}$ once, the composite vertex covering a definite part of it.
  Another possibility could be that the composite vertex the whole region with a constant weight, and we have a construction similar to generalized sections from~\cite{Sen:2015:OffshellAmplitudesSuperstring,deLacroix:2017:ClosedSuperstringField}.

  The first statements concern the poles of the composite vertex at a fixed level of truncation.
The plumbing lattice of a channel is the set of poles $P^2 = - \tfrac{2}{\alpha'} (N + \bar N - 2)$ of the plumbing fixture.
  It is twice as dense as the mass shells of the string field since, over a sector of twist angles narrower than a full turn, the intermediate states need not be level matched.
If the composite vertex has a pole in the pair momentum away from the plumbing lattice of the channels of the amplitude, which we have neither established nor excluded, it is not a compact sub-region with momentum-independent boundary and the standard measure.
Nor is it then a momentum-dependent region of sector type with bounded multiplicity, depending real-analytically on the momenta.
  By this we mean a region, possibly signed, whose points are counted at most a fixed number of times, and which meets the collar of each degeneration in a sector, the set $r < \epsilon(\theta)$ over a range of angles in the plumbing parameter $q = r \e^{\I \theta}$ of the collar, with a width $\epsilon$ bounded away from zero, as the Feynman regions do.
If it has a pole off the mass shells, on the plumbing lattice or not, or no pole at all, it is not a union of Feynman regions either, modulo the standard factorization theorem.
  We have not settled the case of poles on the mass shells only.

To prove the conjecture, we could read the composite vertex as a family of glued surfaces with a weight, which requires $\sprod{2}{2}^{-1}$ to act as a gluing on the off-shell $\Sigma_2$ legs of the two quintic vertices.
  Within a class of gluings we have shown that it does not, for the explicit section of~\cite{Erbin:2026:SL2CVerticesClosed}.
  The class consists of gluings with a weight on the plumbing parameter that does not depend on the inputs, with the circle map of a fixed half whose coordinate at the special puncture is Möbius, and with a gluing operator that acts trivially on the multiplicity spaces, spanned by quasi-primaries, of the matter $\mathfrak{sl}(2)$ algebras generated by $L_{0, \pm 1}$ and $\bar L_{0, \pm 1}$: the ghost zero modes dual to the conformal Killing vectors of the circle, any ghost dressing, and any change of the cut coordinate by Möbius maps, depending on $q$ or not.
  On pairs of tachyons the quartic form is the gluing of the fixed halves over $Q_V$ with the weight $\abs{\rho(q)}^{h_1 + h_2 + h_3 + h_4}$ of~\eqref{p:eq:ws:recipe-primaries}, which depends on the squares of the relative momenta of the two pairs, while the rest of the integrand depends on them only through their scalar product.
  A gluing of the inverse, summed over levels, in this class and on any class of pair states that represents the form on tachyon pairs and separates the even-$L_0^-$ part of the block $W$, would give the form a dependence on the two squares that factorizes.
  The integral over $Q_V$ with this weight factorizes in this way only if $\ln \abs{\rho}$ is constant on $Q_V$, while it vanishes on the inner seam and equals minus the mismatch~\eqref{p:eq:ws:seam-mismatch} at the outer seam at $\theta = \pi$.
  The obstruction therefore appears already at first order in the expansion of the local coordinates.
  Gluings whose cut coordinate is not Möbius, which is how the coordinates of an off-shell vertex move, are outside the class.
  Whether one of them inverts the form is open, as is a reading of the composite vertex as a chain that does not go through a gluing of the inverse.

  A gluing reading needs, in addition, a property of the quintic vertex: that its local coordinates at the two legs of the $\Sigma_2$ pair, in the frame normalized by the circle of the half, are those of the cubic vertex, so that the state on the cut is that of a half.
  It is the second assumption of the circle model below.
  With the quartic vertex of the explicit section, this property fails at first order, in the mapping radii, on an open subset of each face of $\mathcal V_{0,5}$ that borders a Feynman region made of a cubic vertex, a propagator and a quartic vertex, with a pair of punctures other than the $\Sigma_2$ pair pinching off.
  On such a face a vertex consistent with the cubic exchanges carries the coordinates of the exchange diagram~\cite{Zwiebach:1993:ClosedStringField}, so that at the two legs of the pair the quintic vertex has those of the quartic vertex, the coordinates of the half dressed by the input operators of the recipe above, with $\abs{\rho} \neq 1$ near the outer seam of $Q_V$.
  No datum of the quintic vertex enters, and the failure holds for every section consistent with the cubic exchanges, symmetric in its five punctures or not.

  The conjecture can be tested at real momenta, where the composite vertex is a Feynman diagram.
  On six on-shell tachyons at real momenta, the integral of the standard measure over any momentum-independent region of $\mathcal V_{0,6}$ is, up to a constant common to all regions, the Laplace transform of a positive measure in the invariants $p_i \cdot p_j$.
  We define the composite vertex there as the limit of its level truncations, assuming that the quintic vertex is damped as a Gaussian in the Euclidean relative momentum of the pair, that the truncated blocks of the quartic form on pairs are invertible at real momentum off a discrete set, and that the truncations converge.
  A sub-region implies that the ratio of $\vertex{5 \cdot 5}$ to $\vertex{6}$ lies between $0$ and $\tfrac{5}{6}$, and a constant weight implies that this ratio is constant.
  Both imply that the logarithm of $\abs{\vertex{5 \cdot 5}}$ is convex in the invariants and that its growth along a real direction is bounded by the extrema of $\alpha' \ln \abs{z_{ij}}$ over $\mathcal V_{0,6}$.
  A sub-region also implies that the matrices of the values of $\vertex{5 \cdot 5}$, and of $\dvertex{6}$, at the midpoints of finitely many points are semi-definite, with the sign of $\vertex{6}$.
  This is a stronger form of the convexity, which in the slab illustration below separates the cases with a positive measure from the generic one, where the measure is signed.
  These conditions are necessary and not sufficient: a failure refutes the conjecture once the distance between the truncation and its limit is bounded below the margin of the failure, and a pass at sampled points establishes the inequalities at those points only.
  If $\mathcal V_{0,6}$ is invariant under complex conjugation of the moduli, a region holding one point of each conjugate pair carries half of $\vertex{6}$ on tachyons, so that a constant weight $\tfrac{5}{12}$ on the composite vertex is also a covering of the conjectured form.
  The tests need a quintic vertex with its momentum dependence on tachyons.
  In the family of~\cite{Erbin:2026:SL2CVerticesClosed} the region $\mathcal V_{0,5}$ is fixed as a set by the lower vertices, but the vertex is not constructed in its interior, and no local coordinate at a leg is known.

A cusp-shaped region, pinched at a boundary point, can carry a simple pole off the plumbing lattice, and it is a case, simple to draw, of what remains: a momentum-weighted density on a region rather than a region.
We do not know whether a region of this kind reproduces the composite vertex.

The circle picture gives a model of the generic situation, at the price of a second assumption: that the $\Sigma_2$ pair of each quintic vertex is again a rigid half carrying the same circle state.
This is a toy, with no on-shell exactness behind it, and, as shown above, it fails off-shell for every quintic vertex consistent with the cubic exchanges.
We also assume that each quintic vertex is cut along such a circle with a single-valued neck region, which depends on the moduli of the three other punctures; this parametrization is not checked, and the region is not computed.
Under these assumptions, in the model at $\lambda = 1 / 2$ and under the hypothesis on the circle states, the composite vertex is a signed density on the surfaces cut along a circle separating three punctures from three.
In the level-matched sector, for a constant cutoff and quintic neck regions independent of the twist, its measure in the length of the final neck is the derivative of the convolution of the two quintic neck regions.
This derivative is replicated at the shifts $n \tau$ of the tower~\eqref{p:eq:ws:tower}.

For two slabs of widths $\tau_5, \tau_5'$, an illustration, three cases occur.
The measure is a single slab if and only if one width equals $\tau$, and it is then the slab of the other width.
It has compact support if and only if one width is a multiple of $\tau$.
Otherwise it carries signed weight on arbitrarily long necks, with alternating signs, reaching into the Feynman region.
Nothing places a width of the quintic region, which depends on the cluster moduli, on the lattice of a cutoff that is not even constant, so the generic case is the last.
In the first case the composite vertex is a bounded plumbing region in the final neck; in the other two it is neither a union of Feynman regions nor a sub-region with the standard measure, and in the generic case the reduction removes from the sextic vertex a signed density of long-neck surfaces.

At $\lambda = 0$, where the circle states do not fill the block, the exact statement is the projected one.
The sector formula applies to the component of the five-point state that lies in the span of the circle states, and the composite vertex is the sector formula minus a correction carried by the projector,
\begin{equation}
  \label{p:eq:ws:projected-composite}
  \vertex{5 \cdot 5}
    = \vertex{5 \cdot 5}^{\mathrm{sector}} - \vertex{5 \cdot 5}^{\perp},
\end{equation}
where the first term is the signed density above and $\vertex{5 \cdot 5}^{\perp}$ is the gluing form evaluated on the component of the five-point state that the circle states do not reach, the split being exact where both terms converge.
The correction vanishes if the five-point state lies in the span of the circle states, which requires it to obey the selection rule~\eqref{p:eq:ws:selection-rule}.
It is not expected to: the Killing vector behind the rule does not vanish at the three string-field punctures, and already three tachyons give a non-zero component before the reduction by the quintic region and by the inverse of the transform.
That the reduction does not cancel it is argued, not computed; the correction is then expected to be non-zero, nothing in the model making it vanish.
Weakening the hypothesis to the statement that the circle states reach everything the selection rule allows, which the counts through level three support exactly, is therefore not expected to restore the sector formula, since the five-point side is not expected to obey such a rule.
The size of the correction and the position of the poles of the projected composite vertex are open.
The zeros of the compressed determinant move with the compression, as said above, and need not sit at the real zeros of the transform.

The poles of the model at real timelike channel momentum, where their residue does not vanish, are carried by $\dvertex{6}$ and invisible in the amplitude, since $\dvertex{6} + \tfrac{6}{5} \vertex{5 \cdot 5} = \vertex{6}$ is free of them.
All this concerns the untruncated model, and it is not evidence for a pole of the level-truncated composite vertex at the actual vertices.
The position of the density relative to $\mathcal V_{0,6}$, a region defined with the sextic vertex's own coordinates, is not determined by the model.
Whether the coordinate dressing of $\sprod{2}{2}$ projects out the sectors with $L_0^- \neq 0$ of the $\Sigma_2$ line, and with them the real zeros, is open.
Within the model the fate of the real zeros has a precise criterion, whether the circle states contain a vector of a single sector at the zero, but we have not decided it.
The conjecture therefore stays open.
A proof needs a structural argument, a reading of the composite vertex as a chain or a closed form of it, or, if the composite vertex is real-analytic on an open set of invariants, the exact inverse Laplace transform of $\tfrac{6}{5} \vertex{5 \cdot 5}$ in the invariants, compared with the indicator of a region symmetrized under complex conjugation.

\subparagraph{The level-zero kernel.}

The level-zero kernel~\eqref{p:eq:ws:level-0-kernel} is real and symmetric.
For the vertex family of~\cite{Erbin:2026:SL2CVerticesClosed} at the sections of the family considered there with $\beta > \sqrt{2}$, we find that it is moreover strictly positive definite up to an overall sign on the even functions of the relative momentum $q$ at every real value of $P^2$.
The saddle~\eqref{p:eq:ws:level-0-saddle} then has no null function at real $P^2$, and the eigenfunction expansion of the kernel converges absolutely, the operator being trace class.
At complex $P^2$, where a pole of the composite vertex off the plumbing lattice would sit if there is one, the determinant of the discretized kernel vanishes, at $\beta = 2$, at a point off the plumbing lattice.
The zero is found numerically and is stable under the discretization.
We do not know whether the kernel has a null function there, so that, at level zero, a pole off the plumbing lattice is neither realized nor excluded.

\subsubsection{Outlook on a single-strand route}
\label{p:sec:ws:single-strand}



The halves of the circle picture suggest a different cubic theory, with one asymmetric cubic vertex $\vertex{3}'$, with two ordinary punctures and one special puncture, and an auxiliary field with a single string leg.
Such a theory would obtain the quartic vertex by exchanging the auxiliary field between two copies of $\vertex{3}'$ after a change of its kinetic term, as the open-string stub does~\cite{Erbin:2024:OpenStringStub}, and we close with what is established about this route and what is not.
At tree level and four points, ghost and moduli counting allows the quartic vertex to factor through an auxiliary channel in this way only with a rigid cubic vertex and either a tube of bounded length on each strand or the two strands merged into one string leg carrying a single tube.
With two strands on the line the sewing produces four-punctured tori with the genus-one antighost measure rather than $\mathcal V_{0,4}$.
With Möbius discs for the local coordinates the quotient collar of that fixture is numerically not $\mathcal V_{0,4}$ at the parameters sampled.
With a single strand and the reversed plumbing range every count passes with sphere topology, provided the local discs of $\vertex{3}'$ extend beyond the unit circle.
This fails for tangent discs, those of the Sonoda–Zwiebach type as well as those of the cubic vertex of the theory at $\beta = 3 / 2$, and holds for the halves above when the discs at their ordinary punctures are small enough.

Gluing two copies of $\vertex{3}'$ at $\abs{q} = 1$ with the twist alone reproduces the boundary of $\mathcal V_{0,4}$ and not its interior.
At the parameters sampled, the plumbing family covers the interior exactly once when the stub length varies along a seam symmetric under the permutations of the three channels, the seam family, and never for a constant stub length.
This is numerical, and the seam is under-determined.
On-shell, under the condition on the discs above and provided the sewing converges over the reversed plumbing range, the corresponding operator identity reduces to this covering, modulo a measure identity argued by hand.
We have checked the identity numerically at two generic momenta.
The local coordinates transported across the seam agree to first order but not to second, the discs differing by about $15 \, \%$ of their radius.
The strictly cubic single-strand theory is therefore expected, though not shown, to fail gauge invariance off-shell at second order in the seam data on non-primary external states.
A quartic term supported on the collar of the seam would restore it under an embedding condition that we have not checked.

The circle picture also gives the relation between the two formulations.
Up to $c_0^- c_0^+$ in place of $Q_B$, the tower~\eqref{p:eq:ws:tower} that is the inverse of the $2$-string form is the kinetic operator of a single-string auxiliary field glued over the reversed plumbing range, and the gluing region is its propagator.
Neither a cubic theory of this kind with a geometric auxiliary channel, nor the statement that its strictly cubic version is an on-shell coincidence, is established.
In the theory, the cubic vertex fixes the parameter $\lambda$ of the halves at $\lambda = 1 / 2$.
The circle picture is exact on-shell for every $\lambda$, and off-shell it remains a model, in which rigid halves stand in for local coordinates that move with the modulus.
Off-shell, $\lambda$ decides whether the circle states obey the selection rule~\eqref{p:eq:ws:selection-rule}.


\subsection{The vertex region as a dressed gluing, off-shell}
\label{sec:wscircle:offshell}

\status{mixed}\record{113-offshell-dressed-gluing (R93, annotation)}

The block above presents the circle picture as a model, exact on-shell modulo a measure identity, with the dressing of the antighosts~\eqref{p:eq:ws:dressing} and the recipe as statements about how the actual form departs from it.
The derivation in this section turns the recipe into a theorem off-shell: on its vertex region the quartic form, for any quartic vertex given as a section of local coordinates, is the dressed gluing of two copies of a fixed half over $Q_V$, and the rigid gluing of the model is its value on-shell, where the dressing is trivial.
The measure identity, the one step of the block that was argued, follows in two ways, and the Sonoda--Zwiebach sewing of the block follows from it.
What does not follow is a gluing form of the inverse $\sprod{2}{2}^{-1}$: the last part of the section collects what is established about it.

\subsubsection{Assumptions}

We work at genus zero with four external punctures and finite-level inputs in $\Hs_0$, finite linear combinations of Fock states at fixed momenta, with the conventions of \cref{ch:conventions}.
The halves are those of~\eqref{p:eq:ws:half}, glued by $w_1 w_2 = q$.
In the coordinate $u$ of the circle state the punctures of the first copy sit at $u_{1,2} = \pm u_0$, $u_0 = 1 / (2 \beta)$, and those of the second at $u_{3,4} = \pm 2 \beta / q$, so that the cross-ratio of the glued sphere is
\begin{equation}
  \label{eq:wscircle:Zq}
  Z(q)
    = \frac{(q + 4 \beta^2)^2}{16 \beta^2 q}.
\end{equation}
It maps the disc $\abs{q} < 4 \beta^2$ bijectively onto $\mathcal M_{0,4} \setminus [0, 1]$, the closed unit disc onto the $s$-channel Feynman region, and $Q_V$ onto $\mathcal V_{0,4}$ up to a set of measure zero.
The star-shaped form of $Q_V$, which is numerical, is not used: every identity below holds pointwise on $\abs{q} < 4 \beta^2$, and $Q_V$ enters only as the pre-image of $\mathcal V_{0,4}$.

The quartic vertex is given as a section: at each $Z \in \mathcal V_{0,4}$ four local coordinates $f_i(\cdot \,; Z)$, continuously differentiable in $Z$ on each tile of the vertex region, so that the derivatives used below exist.
For the vertex of the explicit section of~\cite{Erbin:2026:SL2CVerticesClosed} this regularity is the hypothesis that the section is real-analytic on each of the twelve tiles, with kinks across the six arcs where tiles meet, which a two-form integral does not see; this is the one hypothesis in force for the theorem applied to that vertex.
The form of the vertex is the standard one: for two tangent vectors $V_1$, $V_2$ of the family,
\begin{equation}
  \label{eq:wscircle:Omega}
  \Omega(V_1, V_2)
    = \kappa \, \big\langle \Sigma \big| \, B[V_1] \, B[V_2] \, \prod_i f_i \circ A_i(0) \big\rangle,
\end{equation}
with $B[V]$ the antighost insertion of a Schiffer vector field representing $V$ and $\kappa$ the normalization of~\eqref{p:eq:ws:p-circ}, fixed by the requirement that the $s$-channel Feynman region reproduce $\int \Omega$ there.
We use the transformation law of a vertex operator under a change of local coordinate, $F \circ A(0) = U_F \, A(0) \, U_F^{-1}$ with $U_{f \circ g} = U_f U_g$, and the two standard facts of the Schiffer formalism: the variation of a surface state is linear in the vector field, its antighost version is the stress-tensor version with $b$ in place of $T$, and the insertion of a vector field depends only on its class modulo the fields that extend holomorphically to the surface minus the punctures.

We write $\langle \Sigma(q); F |$ for the surface state of the glued sphere at $q$ with local coordinates $F = (F_1, \dots, F_4)$ at the external punctures, so that $\langle \Sigma(q); F | A_1 \otimes \cdots \otimes A_4 \rangle = \langle \prod_i F_i \circ A_i(0) \rangle_{\Sigma(q)}$, and $\varphi^h_{2,1}$ for the circle map of a half whose coordinates at the ordinary punctures are $h = (h_1, h_2)$, the same for both copies.

\subsubsection{The quartic form as a dressed gluing}
\label{sec:wscircle:theorem}

\paragraph{The vertex region in the coordinate \texorpdfstring{$q$}{q}.}
Restricted to $Q_V$, the map $Z$ is holomorphic and a bijection onto $\mathcal V_{0,4}$ up to a null set, hence orientation preserving, so
\begin{equation}
  \label{eq:wscircle:pullback}
  \int_{\mathcal V_{0,4}} \Omega
    = \int_{Q_V} Z^* \Omega
    = \int_{Q_V} \dd s \, \dd\theta \; \Omega \big( Z_* \pd_s, Z_* \pd_\theta \big).
\end{equation}
The orientation needs a word because $Q_V$ lies on the reversed plumbing range $s < 0$; reversing the range changes neither the integrand nor the orientation.
The family of pointed surfaces with coordinates $Z \mapsto (\Sigma_Z, f_i(\cdot \,; Z))$, pulled back along $Z(q)$, is the family $q \mapsto (\Sigma(q), F_i(\cdot \,; q))$ with the transported germs
\begin{equation}
  \label{eq:wscircle:transport}
  F_i(\cdot \,; q)
    = N_q^{-1} \circ f_i \big( \cdot \,; Z(q) \big),
\end{equation}
where $N_q$ is the Möbius map of the glued sphere onto the sphere of the section that sends the four punctures to $(0, 1, \infty, Z(q))$.
The two pointed surfaces with coordinates are the same point of the moduli space with local coordinates, so they have the same surface state, $\langle \Sigma_{Z(q)}; f | = \langle \Sigma(q); F(q) |$, and the tangent vectors $Z_* \pd_s$, $Z_* \pd_\theta$ are the tangent vectors $\pd_s$, $\pd_\theta$ of the pulled-back family.
By~\eqref{eq:wscircle:Omega},
\begin{equation}
  \label{eq:wscircle:Omega-q}
  \Omega \big( Z_* \pd_s, Z_* \pd_\theta \big)
    = \kappa \, \big\langle \Sigma(q); F(q) \big| \, B[\pd_s] \, B[\pd_\theta] \, \big| A_1 \otimes \cdots \otimes A_4 \big\rangle,
\end{equation}
with $B[\pd_s]$ the antighost insertion of any Schiffer vector field that represents $\pd_s$ for the family $q \mapsto \langle \Sigma(q); F(q) |$.

\paragraph{A Schiffer vector of the family.}
Split each germ as in step 2 of the recipe, $F_i(q) = h_i \circ g_i(q)$, with $g_i(q)$ a germ fixing $0$; the split is exact for every admissible choice of the $h_i$.
By the transformation law, and by the sewing of the two halves through the plumbing $w_1 w_2 = q$,
\begin{align}
  \label{eq:wscircle:split}
  \big\langle \Sigma(q); F(q) \big| A_1 \otimes \cdots \otimes A_4 \big\rangle
    &
    = \big\langle \Sigma(q); h \big| \, U_{g_1(q)} A_1 \otimes \cdots \otimes U_{g_4(q)} A_4 \big\rangle
  \\
    &
    = \big\langle \varphi^h_{2,1} \big( U_{g_3} A_3, U_{g_4} A_4 \big) \big| \, \e^{- s L_0^+} \e^{\I \theta L_0^-} \, \big| \varphi^h_{2,1} \big( U_{g_1} A_1, U_{g_2} A_2 \big) \big\rangle,
  \nonumber
\end{align}
the last form valid as a level sum for $\abs{q} < 4 \beta^2$, where a circle separating the two pairs exists.
The derivative in $s$ of~\eqref{eq:wscircle:split} has, by the Leibniz rule, one term from the plumbing operator and one from each input operator:
\begin{equation}
  \label{eq:wscircle:leibniz}
  \pd_s \, \e^{- s L_0^+}
    = - L_0^+ \, \e^{- s L_0^+},
  \qquad
  \pd_s U_{g_i(q)}
    = U_{g_i(q)} \, T^{(i)}[v_i^s],
  \qquad
  v_i^s(W)
    = \frac{\pd_s g_i(W)}{g_i'(W)}.
\end{equation}
The first is the plumbing Schiffer vector $- w \pd_w$ at the neck; the second follows from $g(s)^{-1} \circ g(s + \dd s) = \id + \dd s \, v^s$ and the composition law of the $U$'s, with $T^{(i)}[v]$ the stress-tensor contour of $v$ in the frame of input $i$.
Each term is the stress-tensor insertion of a vector field, at the neck in the special coordinate and in disc $i$ in the frame of the input, so the field
\begin{equation}
  v(\pd_s)
    := (- w \pd_w)_{\mathrm{neck}} \oplus \bigoplus_{i = 1}^4 v_i^s,
\end{equation}
defined on the sewing annulus and the four coordinate discs, has $\pd_s$ of the surface state as its stress-tensor insertion.
The Schiffer map from vector fields to tangent vectors of the moduli space with local coordinates is onto, with kernel the fields that extend holomorphically to the surface minus the punctures, the same space that the conservation laws of $T$ and $b$ annihilate; a field whose stress-tensor insertion is $\pd_s$ therefore represents $\pd_s$, and its class is unique.
Replacing $T$ by $b$, and doing the same for $\theta$,
\begin{equation}
  \label{eq:wscircle:dressed-b}
  B[\pd_s]
    = - b_0^+ + \sum_{i = 1}^4 b^{(i)}[v_i^s],
  \qquad
  B[\pd_\theta]
    = \I \, b_0^- + \sum_{i = 1}^4 b^{(i)}[v_i^\theta],
\end{equation}
with the zero modes acting on the neck and $b^{(i)}[v_i]$ acting on $A_i$ before $U_{g_i}$, since the stress-tensor term in~\eqref{eq:wscircle:leibniz} stands to the right of $U_{g_i}$.
Since $v_i(0) = 0$, $b^{(i)}[v_i] = \sum_{k \geq 0} v_{i, k} \, b_k^{(i)}$ contains annihilation-type modes only.
This is the dressing~\eqref{p:eq:ws:dressing}, now derived rather than motivated: with rigid halves the second term of~\eqref{eq:wscircle:leibniz} is absent, and the $q$-dependence of the actual circle states (\cref{sec:wscircle:rigid}) shows that it is not absent for the explicit section.

\paragraph{The dressed gluing over \texorpdfstring{$Q_V$}{QV}.}
Inserting~\eqref{eq:wscircle:split} and~\eqref{eq:wscircle:dressed-b} into~\eqref{eq:wscircle:Omega-q} and~\eqref{eq:wscircle:pullback}, the contribution of the vertex region to the quartic form is
\begin{multline}
  \label{eq:wscircle:theorem}
  \sprod{2}{2} \big|_{\mathcal V_{0,4}} (A_1 \odot A_2, A_3 \odot A_4)
    = \kappa \int_{Q_V} \dd s \, \dd\theta \;
  \\
    \Big\langle \varphi^h_{2,1} \big( U_{g_3} A_3, U_{g_4} A_4 \big) \Big| \, B[\pd_s] \, B[\pd_\theta] \, \e^{- s L_0^+} \e^{\I \theta L_0^-} \, \Big| \varphi^h_{2,1} \big( U_{g_1} A_1, U_{g_2} A_2 \big) \Big\rangle,
\end{multline}
with the disc parts of the antighosts acting on the inputs inside the circle states, the whole a graded product in the order written.
This is~\eqref{p:eq:ws:recipe-plumbing} of the recipe.
The constant is not fitted: over $\abs{q} \leq 1$ the same expression with the coordinates of the Feynman diagram ($g_i = \id$, no disc parts) is the $s$-channel exchange with the propagator, so $\kappa$ is the normalization of the propagator, and it is the same constant on $Q_V$ because $\Omega$ is one form on $\mathcal M_{0,4}$.
On undressed antighosts $B[\pd_s] B[\pd_\theta] = - \I \, b_0^+ b_0^-$, which gives the normalization $- \I \kappa$ of~\eqref{p:eq:ws:p-circ}.

\paragraph{Hypotheses of the theorem.}
For a quartic vertex given as a section with the regularity above, on finite-level inputs, the vertex region's contribution to $\sprod{2}{2}$ is the dressed gluing~\eqref{eq:wscircle:theorem} of two copies of a fixed half over $Q_V$, for any fixed half: this is proven, and for the vertex of the explicit section it holds under the hypothesis that its section is real-analytic on each tile.
With the canonical half, the coordinates of the cubic vertex of the theory, the two copies are $\vertex{3}$ itself, and for a section that matches the Feynman germs at $\abs{q} = 1$ the input operators are trivial on level-matched inputs at that seam.
This last clause rests on the statement that the coordinates of the quartic vertex on $\abs{q} = 1$ are those of the glued cubic vertices, which is argued and measured numerically to converge at rate $2$ as $\abs{q} \to 1^+$ (\cref{sec:wscircle:coordinates}); the core theorem for an arbitrary fixed half does not use it.
The theorem does not give the inverse (\cref{sec:wscircle:inverse}), nor the convergence of the level sum of the integrated form on the fold arcs off-shell, nor anything about inputs of unbounded level.

The theorem shows that the dressing is forced: it is the second term of the Leibniz rule~\eqref{eq:wscircle:leibniz}, present whenever the coordinates of the vertex move with $q$, and the seam condition with the cubic exchanges forces them to move.

\subsubsection{The measure identity on-shell}
\label{sec:wscircle:measure}

\status{Proven}\record{113-offshell-dressed-gluing (R112)}

For four on-shell primaries $A_i$ and any halves of the family~\eqref{p:eq:ws:half}, \eqref{p:eq:ws:half-ordinary},
\begin{equation}
  \label{eq:wscircle:measure-identity}
  - \I \kappa \int_{Q_V} \dd s \, \dd\theta \; \big\langle \varphi_{2,1}(A_3, A_4) \big| \, b_0^+ b_0^- \, \e^{- s L_0^+} \e^{\I \theta L_0^-} \, \big| \varphi_{2,1}(A_1, A_2) \big\rangle
    = \int_{\mathcal V_{0,4}} \Omega(A_1, \dots, A_4),
\end{equation}
so that $\sprod{2}{2}^{\mathrm{circ}} = \sprod{2}{2}$ with $Q = Q_V$ on four on-shell primaries, for inputs of finite level.
The pointwise form of the identity, used in the second proof, holds for every $\beta$ on the whole disc $0 < \abs{q} < 4 \beta^2$; the integrated form holds for $\beta > (1 + \sqrt 2) / 2$, where the Feynman regions do not overlap and $Q_V$ corresponds to $\mathcal V_{0,4}$.
No hypothesis is in force: on-shell no derivative of the section enters.

\paragraph{First proof, from the dressing on primaries.}
Take $A_i = c_1 \bar c_1 V_i \ket{k_i}$ with $V_i$ a matter primary.
Then $b_k^{(i)} A_i = 0$ for $k \geq 0$, since $\{ b_k, c_1 \} = \delta_{k, -1}$ and $b_k \ket{0} = 0$ for $k \geq - 1$, and $L_n A_i = 0$ for $n \geq 1$, since $[L_n, c_1] = (- 2 n - 1) \, c_{n + 1}$ with $c_{n + 1} \ket{0} = 0$ and $V_i$ is primary.
So the disc parts of~\eqref{eq:wscircle:dressed-b} annihilate the inputs and $U_{g_i} A_i = \abs{\rho_i(q)}^{h_i} A_i$, with $h_i$ the eigenvalue of $L_0^+$ and $\rho_i = F_i'(0; q) / h_i'(0)$ the mapping radius, for any germ, Möbius or not.
The theorem~\eqref{eq:wscircle:theorem} becomes~\eqref{p:eq:ws:recipe-primaries}, and on-shell, $h_i = 0$, the weight is one: this is~\eqref{eq:wscircle:measure-identity} for the half $h$ used.
The left-hand side does not depend on the half: for another choice $h'$, $\varphi^{h'}_{2,1}(A_1, A_2) = \varphi^h_{2,1}(U_{h_1^{-1} \circ h'_1} A_1, U_{h_2^{-1} \circ h'_2} A_2)$, and $U$ of a germ fixing $0$ is the identity on a primary of weight zero, by the same mode identities.
So the identity holds for every half of the family.

\paragraph{Second proof, by continuation from the Feynman disc.}
This route knows nothing about the dressing.
Fix on-shell primaries and the halves of the theory, and set
\begin{equation}
  \omega(q)
    := \big\langle \varphi_{2,1}(A_3, A_4) \big| \, b_0^+ b_0^- \, \e^{- s L_0^+} \e^{\I \theta L_0^-} \, \big| \varphi_{2,1}(A_1, A_2) \big\rangle,
  \qquad
  \varpi(q) \, \dd s \, \dd\theta
    := Z^* \Omega.
\end{equation}
\begin{enumerate}
  \setlength{\itemsep}{0.5ex}

  \item $\omega$ is real-analytic on $D := \{ 0 < \abs{q} < 4 \beta^2 \}$.
    Since $b_0^+ b_0^- = - 2 \, b_0 \bar b_0$, it is, up to that factor, the sphere correlator $\langle \prod_i (c \bar c V_i)(u_i, \bar u_i) \; b_0 \bar b_0 \rangle$, with $b_0 = \oint u \, b(u) \, \dd u / 2 \pi \I$ on a circle $\abs{u} = r$ with $u_0 < r < 2 \beta / \abs{q}$, which exists if and only if $\abs{q} < 4 \beta^2$.
    It is evaluated by the residues at $u_1$, $u_2$ of a rational function of the $u_i$, times $\prod_{i < j} \abs{u_{ij}}^{\alpha' k_i \cdot k_j}$, and the four points are distinct on $D$ ($u_3 = u_1$ if and only if $q = 4 \beta^2$, $u_3 = u_2$ if and only if $q = - 4 \beta^2$), so every factor is real-analytic in $(q, \bar q)$ there.

  \item $\varpi$ is real-analytic on $D$.
    On-shell $\Omega = C \, \abs{Z}^{\alpha' k_1 \cdot k_4} \abs{1 - Z}^{\alpha' k_2 \cdot k_4} \, \dd^2 Z$, the Koba--Nielsen form in the frame $(0, 1, \infty, Z)$, independent of the coordinates; $Z$ is holomorphic on $D$ and avoids $0$, $1$ and $\infty$ there (the zeros of $Z$ and $Z - 1$ sit at $q = \mp 4 \beta^2$, the pole at $q = 0$); and
    \begin{equation}
      \label{eq:wscircle:jacobian}
      \dd^2 Z
        = \abs{q Z'(q)}^2 \, \dd s \, \dd\theta,
      \qquad
      Z'(q)
        = \frac{(q + 4 \beta^2)(q - 4 \beta^2)}{16 \beta^2 q^2},
    \end{equation}
    with $Z'$ vanishing only at $q = \pm 4 \beta^2$.

  \item On $0 < \abs{q} \leq 1$, $- \I \kappa \, \omega = \varpi$: this is the statement that the $s$-channel Feynman diagram, two copies of $\vertex{3}$ joined by the propagator $\int \dd s \, \dd\theta \, b_0^+ b_0^- \, \e^{- s L_0^+} \e^{\I \theta L_0^-}$, reproduces the integrand of the amplitude on the $s$-channel Feynman region, the standard factorization~\cite{Erbin:2021:StringFieldTheory}, with the normalization that defines $\kappa$.

  \item The identity theorem for real-analytic functions on the connected open set $D$ gives $- \I \kappa \, \omega = \varpi$ on all of $D$, in particular on $Q_V \subset D$; integrating over $Q_V$ gives~\eqref{eq:wscircle:measure-identity}.
\end{enumerate}
The second proof transports the measure statement of the standard factorization to the reversed range; it does not re-prove it.
It also shows that the identity is not a property of $Q_V$: the pointwise identity holds on the whole disc, and $Q_V$ enters only as the pre-image of $\mathcal V_{0,4}$.

In words, the plumbing measure $\dd s \, \dd\theta \, b_0^+ b_0^-$ is the canonical measure on $\mathcal M_{0,4}$ on the reversed range $s < 0$ as it is on the Feynman disc.
The first proof says why the rigid gluing is exact: the dressing that the actual vertex carries off-shell acts on primaries as a weight, and the weight is one on-shell.
Nothing in this is new to a reader of the general theory of the measure on the moduli space~\cite{Zwiebach:1993:ClosedStringField}: the only non-standard feature is the reversed range, which the continuation handles.

\paragraph{Numerical checks of the identity.}
The pointwise identity of the second proof was checked on four on-shell tachyons with generic exponents: the integrand of the rigid gluing, computed from closed-form ghost and matter correlators with the antighost zero modes of the circle frame as residues at $u_{1,2}$, against $\abs{q Z'(q)}^2 \abs{Z}^{e_{14}} \abs{1 - Z}^{e_{24}}$, at eight points of the Feynman disc and sixteen of $Q_V$, up to $\abs{q} = 8$ at $\beta = 3 / 2$ (a point near the fold arc and one near the outer seam included), for two sets of exponents.
The ratio is constant to $2 \times 10^{-15}$, and equal to one in these normalizations; the ghost correlator used was itself checked against its operator product residues, and the check fails, as it should, with the Jacobians $\abs{Z'}^2$ or $\abs{q Z'}$ in place of $\abs{q Z'}^2$.
The check compares absolute values, so it does not see the sign and orientation of the measure on the reversed range, which rest on the identity theorem of the second proof, and it tests the constancy of the ratio between the disc and $Q_V$ (step 4), not the value of $\kappa$ (step 3).
The mode identities of the first proof were checked in a truncation of the ghost Fock space, where they fail, as they should, on a ghost input that is not primary.

\subsubsection{The inverse of the quartic form off-shell}
\label{sec:wscircle:inverse}

\status{mixed}\record{113-offshell-dressed-gluing (R80 and R93 annotations)}

The theorem writes $\sprod{2}{2}$ on its vertex region as a gluing; it says little about $\sprod{2}{2}^{-1}$, the line of $\Sigma_2$.
We take the inverse on the quotient $\Hn{2} / \ker \sprod{2}{2}$ at fixed total momentum, as the paper does.

\paragraph{Ghost number and the factorization identity.}
The dressing keeps the count of two antighosts, so $\sprod{2}{2}$ has ghost number $+ 2$ and its inverse $- 2$.
By the no-kernel theorem of \cref{ch:worldsheet-dictionary}, no connected kernel of bounded type with the standard sewing insertions equals the inverse: such a gluing over moduli carries ghost number at least $+ 2$.
Dressings with explicit extra antighost contours, each lowering the ghost number by one, can reach $- 2$; they are outside that class and are excluded only as constructions uniform in the vertex data, and a dressing tuned to the actual vertex is not excluded.
What the gluing picture gives positively is a factorization identity.
Let $\hat\varphi$ send $\Sigma_2$ to the function on $Q_V$ whose value at $q$ is the four-component vector of the circle states of the fixed half on the inputs dressed by the four terms of $B[\pd_s] B[\pd_\theta]$, expanded in neck and disc parts, and let $\hat G$ be the operator that pairs through $\e^{- s L_0^+} \e^{\I \theta L_0^-}$ with the corresponding neck insertion and integrates over $Q_V$.
Then~\eqref{eq:wscircle:theorem} reads $\sprod{2}{2} = \transp{\hat\varphi} \, \hat G \, \hat\varphi$, and with $\hat\Pi := \hat\varphi \, \sprod{2}{2}^{-1} \transp{\hat\varphi}$,
\begin{equation}
  \label{eq:wscircle:GPiG}
  \hat G \hat\Pi \hat G \, \hat\varphi
    = \hat G \hat\varphi \; \sprod{2}{2}^{-1} \, \big( \transp{\hat\varphi} \hat G \hat\varphi \big)
    = \hat G \hat\varphi \; \sprod{2}{2}^{-1} \sprod{2}{2}
    = \hat G \hat\varphi,
\end{equation}
that is $\hat G \hat\Pi \hat G = \hat G$ on the image of $\hat\varphi$, for any factorization and with no hypothesis.
This is $\vertex{4} \, \sprod{2}{2}^{-1} \, \vertex{4} = \vertex{4}$ written between circle states: the $c$'s of the inverse cancel the neck antighosts and the moduli of one neighbor.
It says nothing about the form of $\sprod{2}{2}^{-1}$; the lift is bookkeeping, whose only content is that an integral of a sewn correlator factors through a linear map and a symmetric operator.

\paragraph{Statements of the fixed-halves model.}
The model has one circle map independent of $q$ and a gluing operator diagonal in the eigenbasis of $(L_0^+, L_0^-)$, which is what gives the explicit inverse~\eqref{p:eq:ws:inverse-circle}, the tower of necks and the real zeros of $\hat f_{m \neq 0}$.
By the theorem, the actual form differs from the model off-shell exactly by the dressing, and each of the three features fails to transfer for a stated reason: there is no single circle map, the circle states depending on $q$ (\cref{sec:wscircle:rigid}); the dressed antighosts are not functions of $L_0^\pm$, so there is no block form; and the projector clause is a statement about the extended circle map of the model, while for the actual map only the pointwise rank is known, which is necessary and not sufficient.
On-shell the model is the quartic form, by the measure identity, so every on-shell statement of the model is a statement about $\sprod{2}{2}$, and every off-shell one is about the model.

\paragraph{Chains of dressed gluings.}
The Neumann series of \cref{ch:worldsheet-dictionary}, $\sprod{2}{2}^{-1} = \sum_k (D^{-1} R)^k D^{-1}$ for the level-diagonal split, is term by term a chain of $k$ dressed gluings~\eqref{eq:wscircle:theorem}, restricted to their level-off-diagonal parts and sewn along two necks of fixed length, with $k + 1$ algebraic factors; its convergence is the hypothesis that the spectral radius of $D^{-1} R$ is below one, supported in one block.
The Schwinger form $\int_0^\infty \dd t \, \e^{- t \sprod{2}{2}}$ also expands in chains of dressed gluings, but the integral over $t$ does not commute with the expansion, each term diverging, so it gives no term-by-term surface representation.
Open: whether the inverse of the dressed gluing has a closed form of another shape, through an input map that depends on $q$ or the inverse of $\hat G$ fiber by fiber on the lifted target, and whether an operator on the neck, uniform in the level, absorbs the dressing.

\paragraph{Scope of the negative statement.}
``No gluing representation of $\sprod{2}{2}^{-1}$ is established'' covers every connected surface of bounded type with standard sewing insertions, at any genus, the class of the Feynman rules of string field theory, and, in the fixed-halves model, every gluing over a region with the standard measure.
Outside it, and not excluded, are the Neumann series (an infinite sum with algebraic factors), a dressing tuned to the actual vertex, and the signed tower of cylinders with $\delta'$ weights of the model.
The weakest relaxation that evades it is an infinite convergent sum of chains with frozen necks and one algebraic reference inverse, which is the Neumann series under its hypothesis.

\subsubsection{The Sonoda--Zwiebach sewing}
\label{sec:wscircle:sz}

\status{Proven}\record{113-offshell-dressed-gluing (R113)}

Sonoda and Zwiebach~\cite{Sonoda:1990:CovariantClosedString} consider a symmetric Hermitian $\group{SL}(2, \C)$ cubic vertex, with parameter $a$, sewn by the plumbing fixture over a domain of the plumbing parameter that contains $0$ and is the same in the three channels, with level-matched external states.
Their theorem is that no such vertex makes the three channels cover $\mathcal M_{0,4}$ once with coordinates that agree up to a phase at the seams, their condition for the BRST invariance of the four-point amplitude.
The proof runs through the square point $Z = - 1$, and its images $\tfrac{1}{2}$ and $2$, where the seams of the three channels meet.
They allow propagating states that are not level matched.

\paragraph{The sewing on-shell.}
By the measure identity and the Feynman disc, on four on-shell primaries at finite level,
\begin{equation}
  \label{eq:wscircle:sz-sewing}
  \int_{\mathcal F_s} \Omega + \int_{\mathcal V_{0,4}} \Omega
    = - \I \kappa \int_{\{ \abs{q} \leq 1 \} \cup Q_V} \dd s \, \dd\theta \; \big\langle \varphi_{2,1}(A_3, A_4) \big| \, b_0^+ b_0^- \, \e^{- s L_0^+} \e^{\I \theta L_0^-} \, \big| \varphi_{2,1}(A_1, A_2) \big\rangle,
\end{equation}
with $\mathcal F_s$ the $s$-channel Feynman region and $\varphi_{2,1}$ the circle map of $\vertex{3}$, which at $\lambda = 1 / 2$, $\gamma' = \beta$ is their symmetric Hermitian vertex with $a = 2 \beta$.
The right-hand side is their sewing of that vertex over the domain $\{ \abs{q} < \e^{\tau_V(\theta)} \}$ in the $s$-channel, the $t$- and $u$-channels keeping their unit discs.
Over $Q_V$ the range of $\theta$ at fixed $s < 0$ is not the full circle, so the twist integral does not project on $L_0^- = 0$: the propagating states are the components of the circle states that are not level matched, the transforms $\hat f_{m \neq 0}$ being non-zero for even $m$.
This is the enlargement they contemplate, realized.

\paragraph{The square point.}
The points of the disc $\abs{q} < 4 \beta^2$ where $Z(q) = - 1$ solve $(q + 4 \beta^2)^2 = - 16 \beta^2 q$, that is $q^2 + 24 \beta^2 q + 16 \beta^4 = 0$, with roots $q = - 4 \beta^2 (3 \mp 2 \sqrt 2)$; the one on the sheet is
\begin{equation}
  \label{eq:wscircle:square-point}
  q_\square
    = - 4 \beta^2 (\sqrt 2 - 1)^2,
  \qquad
  \abs{q_\square}
    = 4 \beta^2 (3 - 2 \sqrt 2).
\end{equation}
It lies on the ray $\theta = \pi$, where the outer radius is exact,
\begin{equation}
  \e^{\tau_V(\pi)}
    = \bigg( \frac{2 \beta (2 \beta - 1)}{2 \beta + 1} \bigg)^2.
\end{equation}
Then $\abs{q_\square} > 1$ if and only if $4 \beta^2 (3 - 2 \sqrt 2) > 1$, that is $\beta > \beta_\circ := (1 + \sqrt 2) / 2$, and $\abs{q_\square} < \e^{\tau_V(\pi)}$ if and only if $(\sqrt 2 - 1)(2 \beta + 1) < 2 \beta - 1$, again $\beta > \beta_\circ$: both seams give the same threshold.
That the two seam values bound the whole of $Q_V$ on this ray, so that ``between the seams'' means ``inside'', does not use the numerical star shape.
On $(- 4 \beta^2, 0)$ the derivative $Z'(q)$ of~\eqref{eq:wscircle:jacobian} is negative, so $Z$ maps the ray monotonically onto $(- \infty, 0)$, and $Q_V \cap \{ \theta = \pi \}$ is the pre-image of $\mathcal V_{0,4} \cap \R_-$, a single interval because each Feynman region meets $\R_-$ in a single interval, the image of the real diameter of a unit disc under a channel map that is real and monotone on it.
Its end points are the $s$-seam $\abs{q} = 1$ and the cross-channel seam at $\e^{\tau_V(\pi)}$.
So for every $\beta > \beta_\circ$ the square point lies strictly inside $Q_V$; at $\beta = \beta_\circ$ it sits on both seams of the ray at once, $\abs{q_\square} = 1 = \e^{\tau_V(\pi)}$, the region pinches there and the Feynman regions touch at the square point, the configuration of Sonoda and Zwiebach; for $\beta < \beta_\circ$ it is inside the Feynman discs, which then overlap.
By the symmetry of the vertex region under permutations of the punctures, the images $\tfrac{1}{2}$ and $2$ behave alike.

\paragraph{The sewing continued off-shell.}
Continued off-shell, the rigid sewing~\eqref{eq:wscircle:sz-sewing} is the fixed-halves model, and the seam mismatch~\eqref{p:eq:ws:seam-mismatch} shows that it is not consistent at the outer seam: the mapping radii of the two descriptions differ by $\ln ((4 \beta^2 - 1) / 4 \beta)$ at $\theta = \pi$ for every $\beta > \beta_\circ$, vanishing only at the pinch.
This is an obstruction of their type, coordinates that do not agree up to a phase at a seam, met at a seam their construction does not have, between the vertex region and a Feynman region, and through which their proof, which needs the square point on a seam between channels, does not pass for $\beta > \beta_\circ$.
The quartic form of the theory is the dressed gluing~\eqref{eq:wscircle:theorem}, whose coordinates move with $q$ and join the Feynman germs at $\abs{q} = 1$ (for the explicit section, under the hypothesis that it is real-analytic on each tile; the join at $\abs{q} = 1$ is argued, \cref{sec:wscircle:coordinates}): the dressing is what the seam condition forces, and the on-shell sewing is the dressed one with the dressing switched off by $h_i = 0$.

The sewing identity~\eqref{eq:wscircle:sz-sewing} and the position of the square point are proven, the latter in closed form; the seam mismatch is exact at $\theta = \pi$ and numerical at the other seam points; the shape of the domain is numerical at $\beta = 3 / 2$.
That the theory evades their theorem through the quartic vertex, and not through the enlarged space of propagating states, is our reading of these facts.
The square-point computation was checked in exact arithmetic, with the two differences $\abs{q_\square} - 1$ and $\e^{\tau_V(\pi)} - \abs{q_\square}$ factoring through $\beta - \beta_\circ$, together with an independent test, at $\beta = 3 / 2$ and $2$, that $Z = - 1$ lies in the vertex region of the explicit section; at $\beta_\circ$ both differences vanish and at $\beta = 1.1$ the point lies outside.

\subsection{Coordinate data of the explicit quartic vertex}
\label{sec:wscircle:coordinates}

\status{mixed}\record{089-mu-table, 091-per-tile, 092-segment-cert (R86)}

The rank paragraph of the block states that, for the explicit section, the second-order coefficients of the coordinates of a pair vanish at no point of the vertex region other than one surface on its boundary, and, under the hypothesis that the section is real-analytic on each tile, at almost no point for any section of that kind.
This section derives both statements and the value of the coefficients on the boundary $\abs{q} = 1$.
All of it concerns the surfaces and their local coordinates, at the level of geometric data; no correlator normalization enters.

\paragraph{The second-order coefficients.}
Transport the coordinates of the explicit section at $Z(q)$ to the circle frame of the first half, as in step 1 of the recipe, and write $a_i = F_i'(0)$ for the first derivative of the germ at puncture $i$ and $R_i = F_i(\infty)$ for the position of its pole, both in the coordinate $u$.
The conformal Killing vector of the pair is $\beta (u^2 - u_0^2) \pd_u$, and in the local frame of the germ it reads $\pm (\text{dilation}) + (\text{coefficient}) \, W^2 \pd_W$, with second-order coefficients
\begin{equation}
  \label{eq:wscircle:mu}
  \mu_1(q)
    = \beta \, a_1(q) \, \frac{R_1(q) - u_2}{R_1(q) - u_1},
  \qquad
  \mu_2(q)
    = - \beta \, a_2(q) \, \frac{R_2(q) - u_1}{R_2(q) - u_2},
\end{equation}
in the phase in which $a_1 > 0$ and $a_2 < 0$, the convention of the rigid halves, $f_2 = - f_1$.
The vector is a pure dilation in the local frame, the situation of the selection rule, if and only if the germ is Möbius with its pole at the partner puncture, $\mu_i = 0$.
The absolute values and the zero set do not depend on the phase.

\paragraph{The boundary value.}
For two glued halves at parameter $\lambda$ and $\gamma' = \beta$, a direct computation gives $\mu_1 = \mu_2 = \lambda / \beta$ exactly.
Over the $s$-channel Feynman region the coordinates of the explicit section are those of two glued copies of the cubic vertex of the theory, $\lambda = 1 / 2$, which the comparison of maps confirms (and refutes for $\lambda = 0$ and $\lambda = \tfrac{1}{4}$), so $\mu_1 = \mu_2 = 1 / (2 \beta)$ there, $\tfrac{1}{3}$ at $\beta = 3 / 2$, to $10^{-15}$.
From inside $Q_V$, at $\abs{q} = 1 + \varepsilon$, numerically at $24$ angles and given the formulas of the explicit section, $\mu_i - 1 / (2 \beta) = O(\varepsilon^2)$ uniformly in $\theta$: the section joins the Feynman value with a continuous first derivative in the radial direction.

\paragraph{The forced zero at \texorpdfstring{$Z = 1 / 2$}{Z = 1/2}.}
On $Q_V$, $\mu_1 = \mu_2$, by the Klein symmetry of the section.
On the tile obtained from the fundamental domain by the element $g = (\pi, \epsilon)$ of the symmetry group (a permutation and possibly a conjugation), $\abs{\mu_{12}(g \cdot Z_0)} = \abs{\mu_{\pi^{-1}(1) \pi^{-1}(2)}(Z_0)}$ with $Z_0$ in the fundamental domain, so the twelve tiles carry three functions.
At $Z = \tfrac{1}{2}$, that is $q = \pm 4 \I \beta^2$ on the fold circle, a corner of four tiles, the involution $\psi(\zeta) = \zeta / (2 \zeta - 1)$ fixes the punctures $1$ and $2$ and exchanges $3$ and $4$; symmetry of the section under the permutations of the punctures then puts the pole of $F_1$ at $\zeta = 1 = \zeta_2$, the partner puncture, and $\mu_{12} = 0$.
The zero is forced for every section symmetric in its four punctures and continuous there, Möbius or not, by the form of its $2$-jet; it is simple, of slope $2 \sqrt 2 / (3 \beta)$, and its images under the permutations are $\mu_{13} = 0$ at $Z = - 1$ and $\mu_{14} = 0$ at $Z = 2$.
On the closed fundamental domain the explicit section gives $\abs{\mu_{1j}} = \abs{S_0} \, \abs{1 / R - 1 / \zeta_j}$ with $1 / \zeta_j \in \{ 1, 0, 1 / Z \}$, $S_0$ the first-order coefficient of the $s$-channel germ, and closed forms such as $\abs{\mu_{12}(-1)} = 1 / (\sqrt 2 \beta)$ and $\abs{\mu_{1j}(\e^{\I \pi / 3})} = 1 / (\sqrt 3 \beta)$.

\paragraph{The segment criterion.}
On the fundamental domain the explicit section interpolates its pole between two anchors, $1 / R = A + L \, (B - A)$, with a weight $L$ taking values in $[0, 1]$.
A zero of $\mu_{1j}$ therefore needs $t_j := 1 / \zeta_j \in [A, B]$, that is $g_{t_j} := (t_j - A) / (B - A) \in [0, 1]$, and $\mu_{1j} = S_0 (B - A)(L - g_{t_j})$ gives the bound
\begin{equation}
  \abs{\mu_{1j}}
    \geq \abs{S_0} \, \abs{B - A} \; \mathrm{dist} \big( g_{t_j}, [0, 1] \big).
\end{equation}
In the coordinate $x = q / (4 \beta^2)$, with $Z = (1 + x)^2 / (4 x)$ and $\abs{x} < 1$, the anchors are rational:
\begin{equation}
  \label{eq:wscircle:segment-rational}
  B - A
    = \frac{1 + 7 x^2}{6 (1 + x)^2},
  \qquad
  g_1
    = \frac{4 (1 + x^2)}{1 + 7 x^2},
  \qquad
  g_0
    = - \frac{2 (1 + 6 x + x^2)}{1 + 7 x^2},
  \qquad
  g_{1 / Z}(x)
    = g_0(- x).
\end{equation}
The pre-images of $[0, 1]$ under $g_1$, $g_0$ and $g_{1 / Z}$, and the zeros of $B - A$, are explicit and lie outside the closed fundamental domain, except $g_0 = 0$ at $Z = - 1$, where $\mu_{13} = S_0(-1) \, L(-1) \, (1 - 1 / \sqrt 2)$ vanishes if and only if $L(-1) = 0$, as it does for the weight of the explicit section.
Transported to the twelve tiles, $\mu_1 = \mu_2 \neq 0$ on the closed vertex region except at $Z = \tfrac{1}{2}$, and in particular at every point of the open $Q_V$, for every interpolation weight with values in $[0, 1]$ and every $\beta \geq \tfrac{1}{2} + \tfrac{1}{\sqrt 2}$.
This is proven in closed form given the formulas of the explicit section, with no sampling; an earlier version of the argument used a numerical argument principle on a sampled boundary at $\beta = 3 / 2$ and $2$, which agrees with it.
On a grid of $Q_V$ the sampled minima are $\abs{\mu_1} \geq 0.062$ away from the forced zero and $\abs{\mu_1} / \abs{Z - \tfrac{1}{2}} \geq 0.615$ near it, at $\beta = 3 / 2$.

\paragraph{Almost every point, for any section of the kind.}
The statement at every point uses the formulas of the explicit section.
Under the weaker hypothesis that a section is real-analytic on each tile and joins the channel maps continuously at the boundary of $\mathcal V_{0,4}$, together with its symmetry under permutations and complex conjugation and the parameter point $\lambda = \tfrac{1}{2}$, $\gamma' = \beta$ of the cubic vertex, the zero set of $\mu_i$ is of measure zero.
The argument has four steps.
The fundamental tile $T_0$ is open and connected and touches the boundary of the $s$-channel region along an arc with non-empty interior, and no other channel boundary.
By the symmetry, the zero sets of $\mu_1$ and $\mu_2$ on the twelve tiles are the images of the zero sets on $T_0$ of six conditions, each affine in the coefficients of the section.
On the $s$-arc the coordinates are those of the glued cubic vertices, whose poles sit at none of the punctures, so each of the six conditions is violated there; by continuity it is violated on an open subset of $T_0$ next to the arc.
A real-analytic function on a connected open set that is non-zero somewhere vanishes on a null set, which gives the statement on $T_0$ and, by transport, on $\mathcal V_{0,4}$ and on $Q_V$.
Zeros are not excluded: the forced zero at $Z = \tfrac{1}{2}$ lies on the boundary of $T_0$ and is consistent with the statement.

\paragraph{Regularity across the arcs.}
At one tested point of the arcs, $Z = \e^{2 \pi \I / 3}$, for the pair $(1, 2)$, $\mu_1$ is continuous with a jump of its first derivative, carried by $\abs{a_1}$, while the factor $\mu_1 / \abs{a_1} = \beta (R_1 - u_2) / (R_1 - u_1)$ that decides its zeros is continuously differentiable with a jump of its second derivative; across the wall $\Re Z = \tfrac{1}{2}$ and at an interior point nothing jumps.
This is what the hypothesis on the section states, tested at one point; the other arc points and pairs were not run.

\subsection{The actual circle states as rigid states}
\label{sec:wscircle:rigid}

\status{mixed}\record{094-ph31-actual (R92)}

The rank paragraph of the block reduces pointwise questions about the circle map of the quartic vertex to the rigid family at a complex parameter.
This section gives the reduction and the rank computation built on it, for the explicit section, at $\beta = 3 / 2$ unless stated.
Everything is given the formulas of the explicit section; no hypothesis of this document is in force.

\paragraph{The circle state at \texorpdfstring{$q$}{q}.}
At $q \in Q_V$ the circle state is the three-point object of the half, with the rigid maps replaced by the transported germs $F_{1,2}(\cdot \,; q)$ of the section in the circle frame,
\begin{equation}
  \big\langle \chi \big| \varphi_{2,1}(A_1, A_2; q) \big\rangle
    = \big\langle (I \circ \chi)(\infty) \; F_1 \circ A_1(0) \; F_2 \circ A_2(0) \big\rangle,
\end{equation}
the holomorphic sector transformed by $F_j$ and the antiholomorphic one by its conjugate, with $I(w) = 1 / w$.
Over the $s$-channel Feynman region it equals the rigid circle state at $\lambda = \tfrac{1}{2}$, $\gamma' = \beta$ on level-matched inputs, to rounding, and from $Q_V$ it joins it continuously as $\abs{q} \to 1^+$.

\paragraph{Möbius data.}
A Möbius germ is fixed by $F(0)$, $F'(0)$ and $F(\infty)$.
The rigid germ of the half at parameter $\lambda$ is, in the circle frame,
\begin{equation}
  h_1^{(\lambda)}(W)
    = \frac{(\lambda + 1) \, W + \gamma'}{2 \beta \, \big( (\lambda - 1) \, W + \gamma' \big)},
  \qquad
  h_2^{(\lambda)}
    = - h_1^{(\lambda)},
\end{equation}
with $h_1^{(\lambda)}(0) = u_0$ and $h_1^{(\lambda) \prime}(0) = 1 / (\beta \gamma')$ for every $\lambda$, while $\lambda \mapsto h_1^{(\lambda)}(\infty) = (\lambda + 1) / (2 \beta (\lambda - 1))$ is a bijection of $\C$ onto the sphere minus $u_0$.
The section fixes $F_1(0) = u_0$ at every $q$, so
\begin{equation}
  \label{eq:wscircle:mobius-data}
  F_1(W; q)
    = h_1^{(\lambda)} \big( \rho \, W \big),
  \qquad
  \lambda(q)
    = \frac{R_1 - u_2}{R_1 - u_1}
    = \frac{2 \beta R_1(q) + 1}{2 \beta R_1(q) - 1},
  \qquad
  \rho(q)
    = \beta \gamma' \, a_1(q),
\end{equation}
exactly, with complex $\lambda$ and $\rho$; $\rho$ is the mapping radius of the recipe.
Comparing with~\eqref{eq:wscircle:mu}, $\lambda = \mu_1 / (\beta a_1)$, so $\lambda(q) \neq 0$ on $Q_V$ exactly where $\mu_1(q) \neq 0$, that is everywhere but on the surface $Z = \tfrac{1}{2}$.

\paragraph{The Klein symmetry.}
The section realizes the Klein group of the four punctures by Möbius maps, and in the circle frame the unique Möbius map that realizes the exchange $(12)(34)$ on $(u_0, - u_0, 2 \beta / q, - 2 \beta / q)$ is $u \mapsto - u$.
Hence, at every $q \in Q_V$, $F_2 = - F_1 = h_2^{(\lambda)}(\rho W)$, with the same $\lambda$ and the same $\rho$: one scale acts on both legs.

\paragraph{The input rescaling.}
By the two previous paragraphs, $F_i \circ A = h_i^{(\lambda)} \circ (U_\rho A)$ with $U_\rho = \rho^{L_0} \bar\rho^{\bar L_0}$ the operator of $W \mapsto \rho W$, so
\begin{equation}
  \label{eq:wscircle:rigid-reduction}
  \varphi_{2,1}(A_1, A_2; q)
    = \varphi^{\mathrm{rig}}_{2,1} \big( U_\rho A_1, U_\rho A_2 \big) \Big|_{\lambda \text{ in the holomorphic}, \ \bar\lambda \text{ in the antiholomorphic sector}},
  \qquad
  U_\rho \big|_{\Hs_0}
    = \abs{\rho}^{L_0^+}.
\end{equation}
On level-matched inputs $U_\rho$ is real, positive and diagonal in the Fock basis, commutes with the ghost zero modes and with $L_0^\pm$, and is invertible level by level, so the image of $\mathrm{Sym}^2 \Hs_0^{\leq L}$ under the level-projected circle map at $q$ is the rigid image at the parameter $\lambda(q)$.
Its momentum part, $\abs{\rho(q)}^{\alpha' (P^2 / 4 + \Delta^2)}$ for a pair with total momentum $P$ and relative momentum $\Delta$, is a non-zero number at fixed $q$, which a statement at several values of $q$ must keep.

\paragraph{A complex parameter.}
The antiholomorphic sector transforms with $\bar F$, so the rigid state at complex $\lambda$ is the construction with $\lambda$ in the holomorphic and $\bar\lambda$ in the antiholomorphic sector, not the holomorphic continuation of the real family.
Define the rank $r_K(\lambda, \tilde\lambda)$ of the level-projected map with independent parameters in the two sectors; the physical family is the real plane $\tilde\lambda = \bar\lambda$, and the exceptional set at output level $K$ is
\begin{equation}
  E_K
    := \big\{ \lambda \in \C : \; r_K(\lambda, \bar\lambda) < \dim W^{\mathrm{even}}_{\leq K} \big\},
\end{equation}
with $W^{\mathrm{even}}_{\leq K}$ the even-$L_0^-$, ghost-number-four part of the block $W$ through level $K$ at fixed total momentum.
Its finiteness is not automatic: the locus where the rank drops is algebraic in $(\lambda, \tilde\lambda)$, and a component invariant under $(\lambda, \tilde\lambda) \mapsto (\bar{\tilde\lambda}, \bar\lambda)$, an involution whose fixed set is the physical plane, would meet that plane in a curve.

\paragraph{Polynomial structure.}
The rigid germ factors as $h_1^{(\lambda)} = h_1^{(0)} \circ m_\mu$ with $m_\mu(W) = W / (1 + \mu W)$ and $\mu = \lambda / \gamma'$, and the same for $h_2^{(\lambda)}$.
The map $m_\mu$ is the time-one flow of $- \mu W^2 \pd_W$, realized on states by an exponential of $L_1$, so the circle state at $(\lambda, \tilde\lambda)$ is the one at $(0, 0)$ composed with exponentials of $L_1$ and $\bar L_1$ on the inputs, and its entries are polynomials in $\lambda$ and $\tilde\lambda$ of bidegree at most the total levels of the inputs in the two sectors.

\paragraph{The rank through output level two.}
At output level $K \leq 2$, input level $L = 2$, two flat matter components and generic total momentum, the block has $73$ directions.
The rank is $73$ at sampled points of $Q_V$ (exact arithmetic over a finite field at a rational point within $10^{-8}$ of $\lambda(q)$, which is rigorous for generic momenta, and floating point at $\lambda(q)$ itself), and $68$ at $\lambda = 0$, the surface $Z = \tfrac{1}{2}$, the deficit of $5$ of the selection rule appearing only in that limit, linearly in $\lambda$.
The rank-dropping locus has no curve component, the origin is an isolated point of it, regular at first order in every direction, and $E_2$ is finite and meets the real and imaginary axes only at $0$.
Hence the pointwise surjectivity onto the block through level two holds near $\abs{q} = 1$, proven, since maximal rank is an open condition and $\lambda \to \tfrac{1}{2}$ there; at almost every point of $Q_V$, verified by randomized exact arithmetic, the argument resting on the finiteness of $E_2$; and at every point of the open $Q_V$, numerically, from a map of the smallest singular value at step $0.01$ over the image $\lambda(Q_V)$, at one complex and two real values of the total momentum, isolated points of $E_2$ between grid points being excluded only up to the margin of that map.

\paragraph{Relation to the inverse of the integrated form.}
The inverse of the integrated form needs a joint statement, not a pointwise one: the lifted map $\Sigma_2 \mapsto (\pi_W \varphi_{2,1}(\Sigma_2; q))_{q \in Q_V}$, into square-integrable functions on $Q_V$ with values in the even block, must have a dense image modulo the kernel of the gluing operator.
Pointwise surjectivity at almost every $q$ and at every level is necessary for that: if it failed on a set $S$ of positive measure, a measurable choice of unit vectors $\chi_q$ orthogonal to the image at $q \in S$ would give the non-zero continuous functional $f \mapsto \int_S \langle \chi_q, \pi_{\leq K} f(q) \rangle \, \dd s \, \dd\theta$ annihilating the lifted image.
It is not sufficient: by~\eqref{eq:wscircle:rigid-reduction} a lifted image vector depends on $q$ only through $(\lambda(q), \abs{\rho(q)})$, and the complex-conjugation symmetry of the section gives $\lambda(\bar q) = \overline{\lambda(q)}$ and $\abs{\rho(\bar q)} = \abs{\rho(q)}$, so on the curve where $\lambda$ is real and $q$ is not every lifted vector takes equal values at $q$ and $\bar q$, and the lifted statement fails in the topology of uniform convergence, whatever the pointwise rank.
In the square-integrable setting that curve has measure zero.
This is the sense in which the block calls the pointwise statement necessary for the inverse of the integrated form, and not sufficient; the derivation of the closed form of the model's inverse does not apply as written to $q$-dependent circle states, and whether a closed form of another shape exists is open.

\subsection{Direct tests of the gluing recipe}
\label{sec:wscircle:tests}

\status{mixed}\record{097-transformation-law-test, 098-dressing-test (R93)}

The last paragraph of the recipe in the block reports two numerical tests against the explicit section; this section says what each compares and what it leaves out.
Both use the canonical half $\lambda = \tfrac{1}{2}$, $\gamma' = \beta$ and the formulas of the explicit section as input, with no fitted constant.

\paragraph{Surface states, without antighosts.}
The first test compares four-point surface states, with ghost numbers that saturate the sphere and no antighost insertion.
At three points of $Q_V$, $\abs{y} = 0.117$ near the inner seam, $0.222$ and $0.356$ in the interior, $y = q / 4 \beta^2$, $\beta = 3 / 2$, on three sets of inputs of level at most two in each chirality that are not primary (matter descendants in two components at complex momentum, ghost oscillators $c_1 c_{-1}$, $b_{-2} c_1$, $c_1 c_0$, a leg of ghost number zero, both chiralities), it checks two things.
First the transformation law: the correlator with the germs of the section equals the correlator with the coordinates of the half on the inputs $U_g A_i$, with $U_g = \rho^{L_0} \bar\rho^{\bar L_0} \e^{- \nu L_1 - \bar\nu \bar L_1}$ built from the mode algebra and one pair $(\lambda, \rho)$ read from $F_1$ for the four legs, to $3 \times 10^{-14}$ in nine cases.
This compares the Virasoro realization of the germ with the realization of the composite coordinate by Neumann coefficients; without the $L_1$ term, with its sign flipped, or with the parameters continued holomorphically into the antiholomorphic sector, it fails.
Then the sewing: the sum over intermediate states of two fixed halves on the inputs $U_g A_i$ converges to the surface state, to $1.6 \times 10^{-7}$ at $\abs{y} = 0.117$ by intermediate level $12$, and at $\abs{y} = 0.222$ to $7 \times 10^{-6}$ for two input sets at level $12$ and to $1.1 \times 10^{-5}$ for the third at level $14$; at $\abs{y} = 0.356$ the error at level $12$ is between $5 \times 10^{-4}$ and $2 \times 10^{-2}$, falling on average by a factor $0.4$ to $0.5$ per level but not monotonically, consistent and not decided at a small tolerance.
This is evidence that the level sum of the fixed halves converges geometrically off the fold arcs on inputs that are not primary, at $\abs{y} \leq 0.36$ only; the fold arcs are not tested.

\paragraph{The dressed density.}
The second test compares the dressed plumbed density,
\begin{equation}
  \big\langle \varphi_{2,1}(U_g A_3, U_g A_4) \big| \, B[\pd_s] \, B[\pd_\theta] \, \e^{- s L_0^+} \e^{\I \theta L_0^-} \, \big| \varphi_{2,1}(U_g A_1, U_g A_2) \big\rangle,
\end{equation}
with the dressings~\eqref{eq:wscircle:dressed-b} of the Möbius section, against the contact density of the explicit vertex at level $(1, 1)$ in~\cite{Erbin:2026:SL2CVerticesClosed}, times the Jacobian $J = 2 \I \abs{q Z'(q)}^2$ of~\eqref{eq:wscircle:jacobian}.
The inputs are the zero-momentum tachyon $T = c_1 \bar c_1 \ket{0}$ and the ghost dilaton $D$, in the classes $T^4$, $D^4$ and the three classes of $T^2 D^2$; on such legs every intermediate state is built from $c$'s alone and the sum over intermediate states terminates exactly at level two.
At $\abs{q} = 1.05$, at $\abs{q} = 3.5$, and near the outer seam ($\abs{q} = 0.99 \, \e^{\tau_V(\pi)}$, $\theta = - 175^\circ$) for $\beta = 3 / 2$, and at $\abs{q} = 4.5$ for $\beta = 2$, the two sides agree to $6 \times 10^{-10}$, the precision of the finite differences used for the derivatives, and the $T^4$ class to $4 \times 10^{-15}$, both sides at the same normalization of the sphere.
The $b_1$, $\bar b_1$ dressing, the neck signs and the $L_1$ part of $U_g$ each carry weight: dropping the dressing, flipping its sign, dropping $L_1$ or flipping $- b_0^+$ makes the comparison fail at every point.
The coefficients of $b_0^{(i)}$ in the dressing are not tested, since $b_0$ annihilates $T$ and $D$; the case of coordinates that are not Möbius, with $L_{n \geq 2}$ in $U_g$, is not tested either.

\subsection{The patch with three punctures}
\label{sec:wscircle:npatch}

\status{mixed}\record{081-npatch-n3 (R90); 090-h1-error}

The atlas paragraph of the block requires, for $n \geq 3$, patches with $n - 2$ internal moduli and calls the requirement a necessary condition on an object not constructed.
At $n = 3$ the candidate patch can be written down, at the level of surfaces, with no local coordinates at the ordinary legs and no off-shell content.

\paragraph{A candidate patch.}
For $m \in \mathcal V_{0,4}$, the patch $\vertex{4}'(m)$ is the surface of the symmetric quartic vertex in the frame $z = 0, 1, m, \infty$, with the leg at $\infty$ made special and given the Möbius coordinate $w = \lambda(m) / (u - c(m))$, of scale $\lambda$ and center $c$.
The scale is fixed by the requirement that the plumbing of two patches by the twist alone, $\abs{q} = 1$, reproduce the propagator boundary of the region of two quartic vertices: $\abs{\lambda}$ is the transport, equivariant under the permutations, of $\beta \abs{m (m - 1)}^{1/2}$ off the fundamental domain, continuous with a kink across three arcs and their conjugates.
The center $c(m)$ is that of the section, not a free constant, and at $n = 2$, $(\lambda, c) = (\beta, \tfrac{1}{2})$.
The remaining freedom is the seam, the phase of $\lambda$ and the choice of the special leg.
The two-patch family glued by $w_L w_R = q$ maps $(m_L, m_R, q)$ to six points on the sphere, three complex parameters for the three of $\mathcal M_{0,6}$; its cross-ratios are invariant under the Möbius group and the seam; at $n = 2$ it reduces to~\eqref{eq:wscircle:Zq}.
These statements are proven at the level of surfaces.

\paragraph{Census of the Feynman regions of \texorpdfstring{$\mathcal M_{0,6}$}{M06}.}
The Feynman regions of $\mathcal M_{0,6}$ are labeled by the $236$ strata of its compactification, the trees with six labeled leaves and vertices of valence at least three: the open stratum, $15$ regions of a quintic and a cubic vertex, and $220$ built from cubic and quartic vertices ($10$ with two quartic vertices, $45$ with a quartic vertex between two cubic ones, $105$ with four cubic vertices, and $60$ chains of a quartic and two cubic vertices); this count is proven.
For the two patches to cover the vertex region of the sphere with six punctures once, necessary conditions, not jointly sufficient, are local injectivity (a non-zero real $6 \times 6$ Jacobian), global injectivity on the vertex range, non-overlap with each family of Feynman regions above, the pieces that need the quintic vertex, a finite cutoff on every ray, and surjectivity.
That the cutoff is finite on every ray is argued from the degeneration: as $\abs{q} \to \infty$ every ray tends to the divisor of the own $(3|3)$ channel and ends in a region of two quartic vertices, a chain region or a region of four cubic vertices, according to whether two, one or none of the component moduli lie in $\mathcal V_{0,4}$.
None of the other pieces is established.

\paragraph{The center of the special coordinate.}
A simple guess for the center is the centroid $(1 + m) / 3$ of the three ordinary punctures.
Against the explicit section, $c(m) - (1 + m) / 3 = L(m) \, D(m)$ on the fundamental domain, with the interpolation weight $L$ and an explicit factor $D$, so the guess is exact on the six arcs, where $L = 0$, and not only at the symmetric points.
Its error reaches $\tfrac{11}{27}$ at $m = - \tfrac{16}{9}$ and $\tfrac{25}{9}$ on the boundary of $\mathcal V_{0,4}$, exactly, and that this is the global maximum is a measurement; the error relative to $\abs{\lambda}$ is the same on every tile, with maximum $0.122$ in the interior.
It moves the limiting moduli of one sampled ray from outside $\mathcal V_{0,4}$ to inside, and over a grid the membership of the limiting moduli in $\mathcal V_{0,4}$ differs on $16 \, \%$ of the vertex region, so the guess is not adequate and the center has to be read from the section.
The section's center is not holomorphic in $m$ ($\abs{\pd_{\bar m} c} / \abs{\pd_m c} = 0.53$ at $m = - 1.3$, a measurement), so local injectivity needs the real $6 \times 6$ Jacobian and not a complex $3 \times 3$ one, and folds can occur inside tiles.
The scale formula for $\abs{\lambda}$ agrees with the transport of the section on all twelve tiles and on the boundary of $\mathcal V_{0,4}$.
These comparisons are numerical, given the formulas of the explicit section.

\subsection{Blocks from the note on the inverse as a circle gluing}
\label{sec:wscircle:note}

\status{mixed}\record{074-p22-inverse-circle-gluing (R80)}

The two blocks below come from the earlier note on the inverse of the quartic form as a circle gluing.
They give, in more detail than the paper block, the block structure of the inverse in the fixed-halves model, the identification of its two $c$'s with the zero modes of the circle, the reciprocal of the region's transform as a tower of necks, its real zeros, and the six-point composite in the circle toy.
Everything in them concerns the fixed-halves model and, at six points, a further toy assumption on the quintic vertex; their statements have the status listed for the model in \cref{sec:wscircle:paper}, and the projector, the reading of zeros as poles and the composite hold under the hypothesis on the circle states, which the note calls ``the hypothesis''.
Where the first block says that the inverse is ``not a gluing over a reversed region'', it means a gluing over a region of the plumbing parameter with the standard measure: the inverse is not such a gluing, nor, by the ghost-number count, a kernel of the standard sewing class, but it is a signed sum of cylinders with $\delta'$ weights, the tower; and the model it inverts is exact on-shell, the measure identity being proven (\cref{sec:wscircle:measure}).

\begin{quote}
  \small
  \textbf{Notation of the note.}
  The note writes $p$ or $p_{2,2}$ for the quartic form $\sprod{2}{2}$ and $p^{\mathrm{circ}}_{2,2}$ for the model $\sprod{2}{2}^{\mathrm{circ}}$; $\Phi_2$ and $\Phi_2^\dagger$ for the circle map $\varphi_{2,1}$ and its transpose $\transp{\varphi_{2,1}}$; $F$ for the diagonal operator $\hat F$ of the transform; $\sigma$ and $\sigma_V(\theta)$ for the cutoff $\tau$ and the outer radius $\tau_V(\theta)$ of $Q_V$; $\mc V_n$ for $\vertex{n}$, $W_6$ for the dressed sextic vertex $\dvertex{6}$, and $\mc C$ for the composite $\vertex{5 \cdot 5}$.
  The block $W = c_0^+ c_0^- \Hs^\downarrow$, the transform $\hat f_m(h)$, the projector $P_{\mc I}$ and the normalization $\kappa$ are those of the paper.
  No reality condition enters these blocks; where the older notes write $\Sigma^\ddagger = \Sigma$, the convention of this document is $\Sigma_k^\star = (-1)^k \Sigma_k$, which differs by a sign for odd $k$.
\end{quote}

\subsection{The inverse}
\label{icg:sec:inverse}

\subparagraph{Block structure and the projector.}
Write $\mc H^\downarrow := \ker b_0^+ \cap \ker b_0^-$ and $W := c_0^+ c_0^- \mc H^\downarrow$, with $c_0^\pm = \tfrac12 (c_0 \pm \bar c_0)$, so that $\mc H = \mc H^\downarrow \oplus c_0^+ \mc H^\downarrow \oplus c_0^- \mc H^\downarrow \oplus W$.
In the eigenbasis $\ket{h, m}$ of $(L_0^+, L_0^-)$ the operator of~\eqref{p:eq:ws:p-circ} is
\begin{gather}
  \label{icg:eq:G}
  G
    := - \I \kappa \, b_0^+ b_0^- \, F,
  \qquad
  F \ket{h, m}
    = \hat f_m(h) \ket{h, m},
  \\
  \label{icg:eq:fhat}
  \hat f_m(h)
    := \hat f_{Q_V}(h, m),
  \qquad
  \hat f_R(h, m)
    := \int_{R} \dd s \, \dd\theta \; \e^{-s h + \I m \theta},
\end{gather}
the Laplace transform in the neck length and Fourier transform in the twist of the region.
$G$ maps $W$ to $\mc H^\downarrow$ and annihilates the other three summands, so the form $\braket{\cdot, G \, \cdot}$ lives on $W \times W$, non-degenerate wherever $\hat f_m(h) \neq 0$; as a form on $\mc H^{\wedge 2}$, $p = \Phi_2^\dagger G \Phi_2$ with $\Phi_2^\dagger$ the transpose of the circle map for the pairing.

\subparagraph{The projector.}
The inverse $p^{-1}$ is defined on $\mc H^{\wedge 2} / \ker p$; pushing it back to $\mc H$ requires a choice of section, and this is where the kernel of $p$ enters.
With such a choice, $\Pi := \Phi_2\, p^{-1} \Phi_2^\dagger$.
From $p = \Phi_2^\dagger G \Phi_2$ one line gives $G \Pi G = G$ on the image $\mc I$ of the circle states, when paired against $\mc I$.
Moreover, $P_{\mc I} := \Pi G$ is the $G$-orthogonal projector onto $\mc I$, modulo $\mc I \cap \mc I^{\perp_G}$.
Since $G$ sees only the $W$-components of the circle states,
\begin{equation}
  \label{icg:eq:Pi}
  \Pi
    = P_{\mc I} \, G_W^{-1},
\end{equation}
with $G_W^{-1}$ the inverse of $G$ on the block, $\mc H^\downarrow \to W$: the inverse of the restricted form is the block inverse followed by a projection, and the projection is the whole difference.
It is invisible when both neighbors are circle states (\cref{p:sec:ws:telescoping}).

\subparagraph{The hypothesis.}
$G P_{\mc I} = G$ on all of $W$ – the projector is invisible to the gluing – only under a hypothesis: that the projection of the circle states onto $W$, extended to $\Sigma_2$ of unbounded level (the space in which the field lives), reaches the whole even-$L_0^-$, ghost-number-4 part of $W$ at fixed total momentum, modulo $\ker G$.
It would be falsified by a persistent rank deficit of the projected circle map, level by level, against the dimension of the even block.
This is motivated – at fixed $P$ the inputs include arbitrary functions of the relative momentum and all oscillator states at both punctures while each level subspace of $W$ is finite-dimensional, and the matter part of the map is triangular in level with an invertible leading part – but the ghost sector is the expected failure mode, since $c(x) c(-x)$ starts at $c \pd c$.
Under the hypothesis $p^{-1}$ acts on $\Sigma_2$ of unbounded level, outside the finite-level setting in which every pairing of circle states is a finite correlator; density of the finite-level image would not suffice, because $\Phi_2^\dagger$ sums over all levels.

\subparagraph{The two $c$'s are the circle's zero modes.}
On $\mc H^\downarrow$, with $\{b_0^\pm, c_0^\pm\} = 1$, $\{b_0^\pm, c_0^\mp\} = 0$ and $b_0^\pm \ket{\downarrow} = 0$,
\begin{equation}
  \label{icg:eq:zero-modes}
  b_0^+ b_0^- \, c_0^- c_0^+ \ket{\downarrow}
    = b_0^+ (1 - c_0^- b_0^-) c_0^+ \ket{\downarrow}
    = \ket{\downarrow},
  \qquad
  c_0^- c_0^+ \, b_0^+ b_0^- \big|_W
    = 1,
\end{equation}
so $(b_0^+ b_0^-)^{-1} = c_0^- c_0^+$ between $W$ and $\mc H^\downarrow$, and in the fixed-halves model
\begin{equation}
  \label{icg:eq:inverse}
  \Phi_2 \, (p^{\rm circ}_{2,2})^{-1} \, \Phi_2^\dagger
    = P_{\mc I} \; \frac{\I}{\kappa} \; c_0^- c_0^+ \; F^{-1},
  \qquad
  F^{-1} \ket{h, m}
    = \frac{1}{\hat f_m(h)} \ket{h, m}.
\end{equation}
The two $c$'s of $p^{-1}$ are the zero modes dual to the two Schiffer antighosts of the gluing, i.e. to the dilation and the rotation of the circle: the residual $\C^*$ of the two-punctured sphere, the pair of conformal Killing directions the counting asks for.
This identification is a fixed-halves statement: the dressed antighosts~\eqref{p:eq:ws:dressing} contain $b_n$ with $n \neq 0$ in the circle frame, and for $p_{2,2}$ itself only the ghost-number count transfers.

\subparagraph{Not a gluing over a reversed region.}
The group law $q^{L_0} \bar q^{\bar L_0} \, q'^{L_0} \bar q'^{\bar L_0} = (qq')^{L_0} (\bar q \bar q')^{\bar L_0}$ holds pointwise, but the $b$-dressed gluings $\mc G_R = \int_R \dd s \, \dd\theta \, b_0^+ b_0^- \e^{-s L_0^+} \e^{\I\theta L_0^-}$ compose to $\mc G_R \mc G_{R'} \propto (b_0^+ b_0^-)^2 = 0$: a gluing absorbs two units of ghost number through its two antighosts, an inverse must supply them, so it carries $c$'s rather than a reversed parameter.
At the level of the modulus, the inverse of an integral over a region is not the integral over the inverse region: for one sector and a constant cutoff $\sigma$,
\begin{equation}
  \label{icg:eq:no-inverse-region}
  \hat f_{[-\sigma, 0]}(h) \; \hat f_{[0, \sigma]}(h)
    = \frac{(\e^{\sigma h} - 1)(1 - \e^{-\sigma h})}{h^2}
    \neq 1 .
\end{equation}
Only a point inverts to a point (a frozen neck to the frozen neck of opposite length), and a two-form has no point to invert.

\subparagraph{The reciprocal of the region's transform.}
For a rotation-independent cutoff, $Q = \{1 < \abs{q} < \e^{\sigma}\}$, the twist integral gives $2\pi \delta_{m,0}$ and the neck integral the transform of the slab $[-\sigma, 0]$,
\begin{equation}
  \label{icg:eq:tower}
  \hat f_m(h)
    = 2\pi \, \delta_{m,0} \, \frac{\e^{\sigma h} - 1}{h},
  \qquad
  \frac{1}{\hat f_0(h)}
    = \frac{1}{2\pi} \, \frac{h}{\e^{\sigma h} - 1}
    = \frac{h}{2\pi} \sum_{n \ge 1} \e^{-n \sigma h}
  \quad (\Re h > 0),
\end{equation}
the other branch, $\Re h < 0$, being the one we call reversed.
Since $h \, \e^{-n\sigma h} = \int \dd s \, \delta'(s - n\sigma) \, \e^{-sh}$, the reciprocal is the Laplace transform of $\sum_{n \ge 1} \delta'(s - n\sigma)$: a discrete tower of necks of positive lengths $n\sigma$, each carrying one derivative in its length.
In operator form each term is $c_0^- c_0^+ L_0^+ \e^{-n\sigma L_0^+}$ with $L_0^+ = \{Q_B, b_0^+\}$, the zero-mode part of the $\Psi$ kinetic operator $c_0^- Q_B$ on a cylinder of length $n\sigma$; the $L_0^+$ is $-\pd_s$ of the plumbing family, a derivative in the neck length acting on whatever is glued to the cylinder.

Up to $c_0^- c_0^+$ in place of $Q_B$ this is the kinetic operator $Q_B L_0^+ / (\e^{\sigma L_0^+} - 1)$ of an auxiliary field with one string leg, obtained by gluing over the reversed plumbing range: between the $\Sigma_2$ formulation and the single-string one the roles of region and tower are exchanged, the tower being the kinetic operator there and the region its propagator.
The same exchange appears for the string field itself: the zero-mode part of the $\Psi$ kinetic form is $c_0^- c_0^+ L_0^+$, the point $s = 0$ with a $\delta'$ (transform $h$), and its inverse is the half-line $[0, \infty)$ (transform $1/h$); the half-line inverts to a single $\delta'$ because it is closed under addition, while a bounded slab inverts to the $\delta'$ at one endpoint convolved with the geometric series generated by the other – that series is the whole difference between $Q_B^{-1}$ and $(p^{\rm circ}_{2,2})^{-1}$ at constant cutoff.

In the level-matched sector the reciprocal~\eqref{icg:eq:tower} is regular at $h = 0$ (value $1/(2\pi\sigma)$): the $\Sigma_2$ line has no mass-shell pole, as for the open-string auxiliary of~\cite{ErbinFirat2023}.
Its poles sit at $h = 2\pi \I k / \sigma$, $k \neq 0$, the zeros of $\hat f_0$: complex momenta, zeros of the block operator $G_W$ on that sector, and poles of $p^{-1}$ only under the hypothesis above (one sector vanishing does not make $\det p$ vanish, every level of a circle state being populated).

The actual cutoff $\sigma_V(\theta)$ is not constant; then $\hat f_m(h)$ is the $m$-th Fourier coefficient of $(\e^{\sigma_V(\theta) h} - 1)/h$, and odd $m$ vanish by $\pi$-periodicity, consistently with the circle states having no odd-$L_0^-$ component.
In the level-matched sector, with $\nu$ the distribution of $\sigma_V(\theta)$ for $\theta$ uniform on $[0, 2\pi)$, a probability measure on $[\ln\tfrac94, \ln 9]$ at $\beta = 3/2$, and $\hat\nu(h) := \int \dd\nu(\sigma) \, \e^{\sigma h}$,
\begin{equation}
  \label{icg:eq:tower-sigmaV}
  \frac{1}{\hat f_0(h)}
    = \frac{1}{2\pi} \, \frac{h}{\hat\nu(h) - 1}
    = - \frac{h}{2\pi} \sum_{n \ge 0} \hat\nu(h)^n
  \quad (\Re h < 0),
\end{equation}
a tower of cylinders whose lengths are sums of $n$ draws from the seam distribution, on the reversed branch only (on the other branch there is a tower only if $\nu$ is a point mass); none of its zeros is real, since $(\hat\nu(h) - 1)/h > 0$ for real $h$.
In the even sectors $m \neq 0$, however, $\hat f_m(h) = \frac1h \int_0^{2\pi} \e^{\I m\theta} \e^{\sigma_V(\theta) h} \, \dd\theta$ is real for real $h$ (the profile is even in $\theta$), and its asymptotic signs are fixed by the geometry: as $h \to +\infty$ the flat fold arcs $(\theta_1, \theta_2)$ dominate, with sign $\operatorname{sgn}[\sin m\theta_2 - \sin m\theta_1]$, positive for $m = 4$, while as $h \to -\infty$ the minima at $\theta = 0, \pi$ dominate and every even $\hat f_m$ is negative.
So $\hat f_4$ has an odd number of real zeros.
Numerically (at $\beta = 3/2$, from the computed profile, three digits, scanning $h \in [-10, 10]$), $\hat f_4$ has one, at $h_* \approx -2.36$, $\hat f_6$ has two, at $-4.89$ and $-1.91$, and $\hat f_2$ has none.
Since $h$ is the $L_0^+$ eigenvalue of the circle state, $h = \tfrac{\alpha'}{2} P^2 + N + \bar N - 2$ with $P$ the momentum through the circle, these are zeros of $G_W$ at \emph{real timelike} momentum through the circle, $\tfrac{\alpha'}{2} (-P^2) = N + \bar N - 2 - h_*$ (at six points, in the channel invariant $s_{123}$ of \cref{icg:sec:six-points}), in sectors $N - \bar N = m \neq 0$ that contain no level-matched state; they are never reached at real Euclidean momentum, where $h \ge N + \bar N - 2$, which is $\ge 2$ in the $m = 4$ sector.
Under the hypothesis they are poles of $p^{-1}$; their fate in the composite is open (\cref{icg:sec:six-points}).




The second block of the note is the six-point composite in the circle toy, which the six-point paragraph of the paper block summarizes.
It rests on two assumptions besides the hypothesis on the circle states: that each quintic vertex is cut along a circle with a single-valued neck region, a parametrization not checked, and that the $\Sigma_2$ pair of each quintic vertex is again a rigid half, a toy assumption with no on-shell exactness behind it, which is known to fail off-shell for a quintic vertex symmetric in its five punctures.
The trichotomy of the neck measure in the widths is proven and checked; the position of the density relative to $\mathcal V_{0,6}$ is argued, not determined by the model.

\subsection{Six points}
\label{icg:sec:six-points}

Each $\mc V_5(\Psi^3, \Sigma_2)$ is cut along a circle enclosing its two $\Sigma_2$ legs: a disc carrying the three $\Psi$'s (two real moduli $m_3$) glued to a circle state over a $q$-region $Q_5(m_3)$ that depends on the cluster moduli (we assume such a parametrization with a single-valued neck region; it is not checked, and $Q_5$ is not computed).
To say anything explicit one has to assume in addition that the $\Sigma_2$ pair of each $\mc V_5$ surface is again the rigid $\mc V_3'$ half with its transported coordinates, so that the circle state is the same $\Phi_2(\Sigma_2)$ as in $p$ and the gluing operator $G_5(m_3) := -\I\kappa\, b_0^+ b_0^- F_5(m_3)$, with $F_5(m_3) := \int_{Q_5(m_3)} \dd s \, \dd\theta \, \e^{-s L_0^+} \e^{\I\theta L_0^-}$, is diagonal in $L_0^\pm$, with eigenvalues $\hat f_5(h, m; m_3) := \hat f_{Q_5(m_3)}(h, m)$.
This is a toy assumption, with no on-shell exactness behind it (the $\Sigma_2$ legs of $\mc V_5$ are always off shell in the composite) and known to fail for an $S_5$-symmetric $\mc V_5$ off shell; with the actual coordinates $G_5$ carries inner antighost dressings and nothing below has a diagonal form.
Under it, and under the hypothesis that the circle states reach the block, the composite is
\begin{gather}
  \label{icg:eq:composite}
  \mc C
    = \mc V_5 \, p^{-1} \, \mc V_5
    = \int \dd m_3 \, \dd m_3' \;
      \Braket{\Phi_3^{[b]}(m_3), \; G_5(m_3) \, \Pi \, G_5(m_3') \, \Phi_3^{[b]}(m_3')},
  \\
  \label{icg:eq:composite-block}
  G_5(m_3) \, \Pi \, G_5(m_3')
    = - \I\kappa \; b_0^+ b_0^- \; F_5(m_3) \, F^{-1} \, F_5(m_3')
  \quad \text{on } W,
\end{gather}
with $\Phi_3^{[b]}$ the three-puncture circle state dressed by the two cluster antighosts, which stay inside the cluster disc.
We draw six consequences.

\begin{enumerate}[leftmargin=*, itemsep=0.3ex]
  \item \emph{A signed density on the $(3|3)$-cut surfaces.}
    Of the eight antighosts of the two $\mc V_5$'s the two $c$'s remove the $b_0^\pm$ of one circle, leaving six, the top degree of $\mc M_{0,6}$; of the eight real parameters $(m_3, q_1, m_3', q_2)$ only six are moduli of the glued surface, $(m_3, m_3', Q)$ with $Q$ the parameter of the $\{123\}|\{456\}$ neck, and the remaining two form a fiber over which the integrand is a scalar weight.
    In the level-matched sector with a constant cutoff and twist-independent $\mu_5$, on the branch $\Re h > 0$, the neck measure in the length $S$ of the final neck is
    \begin{equation}
      \label{icg:eq:neck-measure}
      \mu_{\mc C}(S; m_3, m_3')
        = \sum_{n \ge 1} \pd_S \big( \mu_5(\cdot\,; m_3) * \mu_5(\cdot\,; m_3') \big)(S - n\sigma),
    \end{equation}
    with $\mu_5$ the indicator of $Q_5(m_3)$ in the neck length: the $S$-derivative of the convolution of the two $\mc V_5$ neck regions, replicated at the shifts $n\sigma$; the $c$'s act on the neighbors' moduli as $\delta'$, as in \cref{p:sec:ws:telescoping}.
  \item \emph{The trichotomy in the neck widths.}
    For slabs of widths $\sigma_5, \sigma_5'$ (an illustration: $Q_5$ is not computed and $Q_V$ has no constant cutoff), the measure~\eqref{icg:eq:neck-measure} is a single slab if and only if one width equals $\sigma$, say $\sigma_5 = \sigma$, and it is then the \emph{other} width's slab $[-\sigma_5', 0]$; compactly supported if and only if one width is a multiple $k\sigma$, and then $k$ overlapping slabs on the reversed side; otherwise it carries signed weight on arbitrarily long necks, alternating in sign and without decay at the level of the measure.
    In transform, $\hat f_{\sigma_5} \hat f_{\sigma_5'} / \hat f_\sigma$ with $\hat f_x = (\e^{xh} - 1)/h$ is a single $\hat f_c$ only if the exponent sets $\{\sigma_5 + \sigma_5', \sigma_5, \sigma_5'\}$ and $\{\sigma + c, \sigma, c\}$ agree, and is entire only if every pole $2\pi\I k/\sigma$ is cancelled by a zero of a numerator, which at $k = 1$ already forces one width into $\sigma\N$.
    The sum of the widths is not the criterion: $(\sigma/2, 3\sigma/2)$ leaves weight at arbitrarily large $S$.
    Nothing places a width of $Q_5(m_3)$, which depends on $m_3$, on the lattice of a $Q_V$ that has no constant cutoff, so the generic case is the last: $\mc C$ reaches into the $\{123\}|\{456\}$ Feynman region with long necks of lengths near $n\sigma$.
    In all three cases it is neither a union of Feynman regions nor a sub-region with the standard measure.
  \item \emph{Poles, and where they sit.}
    The function $\hat f_5 \hat f^{-1} \hat f_5'$ is meromorphic in $h$ with poles at the zeros of $\hat f_m$ unless the residue vanishes, i.e. unless the $m_3$-integrated projection of $\mc V_5(\Psi^3, \cdot)$ onto the degenerate sector at that momentum vanishes (a pointwise zero of $\hat f_5(h; m_3)$ moves with $m_3$ and cancels nothing after integration).
    In the level-matched sector these are complex momenta; in the even $m \neq 0$ sectors, for the actual $\sigma_V$, they sit at real timelike channel momentum, $\tfrac{\alpha'}{2} s_{123} = N + \bar N - 2 - h_*$ with $s_{123} := -(k_1 + k_2 + k_3)^2$ the invariant of the momentum through the $\{123\}|\{456\}$ circle, e.g. $N + \bar N + 0.36$ for $m = 4$, at non-integer values where no string mass sits.
  \item \emph{Carried by the dressed sextic vertex, invisible to the amplitude.}
    Since $W_6 + \tfrac65 \mc C = \mc V_6$, the six-point amplitude does not see these singularities; the dressed vertex $W_6$ does and, wherever the residue does not vanish, is not entire in the momenta.
    Geometrically the dressing removes from the sextic vertex a signed density of long-neck surfaces, with the neck profile~\eqref{icg:eq:neck-measure}, and not a sub-region of $\mc V_{0,6}$.
  \item \emph{What this says about which region of $\mc M_{0,6}$ the composite covers.}
    Where the measure is a single slab, $\mc C$ lives on $\abs{Q} \in (1, \e^{\sigma_5'})$, on the vertex side of the boundary $\abs{Q} = 1$ of the $\{123\}|\{456\}$ Feynman region; its position relative to $\mc V_{0,6}$, a region defined with the sextic vertex's own coordinates, is not determined by the model, and the sum over the ten $(3|3)$ channels is not addressed.
    The region arithmetic $R_5 \cdot R^{-1} \cdot R_5'$ fails for three reasons: there is no inverse region (the inverse is a tower of $\delta'$'s), the two $\mc V_5$ neck parameters form a fiber, and a width off the lattice leaves weight on long necks.
  \item \emph{What this does not say about poles of the level-truncated composite.}
    The mechanism of item~3 is realized in the untruncated model, under the surjectivity hypothesis and the toy assumption; it is not evidence for a pole of the level-truncated composite at the actual vertices.
    Whether the actual $p_{2,2}$, through its coordinate dressing, projects the non-level-matched sectors of the $\Sigma_2$ line out – which would remove these sectors, and with them the real zeros, from the $\Sigma_2$ line – is open.
\end{enumerate}



